\chapter{Layers of added complexity}\label{ch:LayAdd}
In order to better understand the languages in the region as they are spoken today,
it helps to also understand the context in which they are spoken.
That is, the geography,
prehistory, and patterns of contact and trade.
It is also important to understand the history of the study of the languages,
and some of the conventions used in earlier eras to document these languages.

In looking to classify the Austronesian languages in the region,
one cannot wear Austronesian blinders.
When people speaking Austronesian languages first entered the region they did not come into a vacuum.
When they first began to arrive in this diverse region spanning almost 2,000 kilometres from west to
\ili{east}, just over 1,000 kilometres north to south,
and encompassing over 2,000 islands,
there were already other people on many of those islands speaking diverse languages that we subsume loosely
under the term Papuan.\footnote{
		Papuan does not define a single genealogical linguistic affiliation.
		It is a cover term of convenience used to refer to numerous unrelated language families in the greater
		New Guinea region that are not Austronesian,
		and not Australian Aboriginal
		(\citealp{Foley1986}, \citealp[1]{Klamer2017}).}

The speakers of Austronesian languages encountered new geography,
new climates, new flora
and fauna, new food possibilities, and new cultural ideas they had not previously had ways of talking
about, nor had they had any need for them.
\ili{As} they migrated, traded,
became bilingual and intermarried with other groups,
they encountered and accommodated new kin terms
and social categories that were deemed useful.
Different groups in different islands had significantly different patterns of contact and trade.
So, for example, some terms
and patterns found along the southern group of islands,
are not found in the northern islands, and vice versa.
And there are more localised patterns as well.
Some languages adopted new words.
Closely related neighbouring languages did not adopt those same words.
Some new words became conventionalised across whole clusters of languages; some only in individual languages.
Examples illustrating these phenomena are given in this and subsequent chapters.

Some early sources documenting languages in the region
were written by people with some training in (European or classical) languages.
Others were compiled by people with no special training.
So it should not be surprising that there is great variation in the notation conventions,
accuracy and reliability of the early documents,
as well as in the details.
Some of the early notational conventions were fairly widespread during a certain time period,
but not necessarily used consistently.
So, there are many
and various ways in the region to represent glottal stops,
phonetically long vowels, vowel systems that have more than five vowels,
word stress, ``funny'' consonant sounds, allophones, word breaks, etc.
Or these sounds were not represented at all by some people documenting the languages.
Some investigators used three different ways to indicate vowel length within a single language.
Some were consistent in their different practices for marking different long vowels,
and some were not.
Some body parts and kin terms are given in first person,
second person or third person forms.
Some investigators give just a noun or verb root,
with an indication that there is an inflection paradigm,
but without providing such a paradigm.
So we have some roots,
but not necessarily in the form of a full word that would be recognised by speakers of the languages.
Similarly for verbs, some investigators provide verbs with their inflections,
and others cite only the root,
without the information necessary to have a citation form,
and without a full paradigm for any verb.
This is complicated in that languages in many localised regions have morphophonemic changes in their
verbal paradigms, and the form provided may not directly reflect the normal historical sound correspondence for that language,
but rather reflects the result of a subsequent morphophonemic change.
European investigators often glossed local flora
and fauna with names of similar European flora
and fauna, and (in the classical tradition) cited verbs as infinitives which were actually inflected
finite verbs in the local languages.

All of these issues add layers of complexity to the task of unravelling the prehistory
and classification of languages in the region.
In \srf{02-01} we look at what is known about the geography,
climate, and flora and fauna in relation to the early prehistory of humans in the region,
providing a backdrop to the later arrival of the Austronesian speaking peoples.
In \srf{02-02} we look at evidence for various kinds of on-going contact,
migrations, and trade that continued after the arrival of the Austronesian speaking peoples,
showing that there were many external factors that continued to impact the development
and ongoing diversification of languages in the region.
In \srf{02-03} we look at various kinds of evidence found in the languages themselves for the kinds of
contact-induced change described in the previous two sections.
And in \srf{02-04} we describe various issues in understanding
and unravelling the transcription of the languages necessary for careful comparative work.
Some of these issues have been ignored or misunderstood by comparative linguists,
and this has contributed to confusion with regard to the languages under investigation in this book.

\scp{Wallacea and its peopling}{02-01}
Before focusing on the prehistory
and classification of the Austronesian languages of the region,
it is helpful to understand the context.
There are two biogeographical regions with relatively shallow seas connecting islands within each of
them, and relatively deep water crossings between the two zones.
\it{Sundaland} extends from mainland South\ili{east} Asia as far \ili{east} as the islands of Java, Bali, and Borneo.
\it{Sahulland} includes Australia, New Guinea,
and the \ili{Aru} Islands in southeastern Maluku.
(\frf{02-01}.)

\begin{figure}
	\includegraphics[width=\columnwidth]{WallaceLine.pdf}
	\fcp{Biogeographical zones with Wallacea as a transition zone}{02-01}
\end{figure}

The \it{Wallace Line} in \frf{02-01} indicates the westernmost extent of fauna (such as marsupials) commonly associated with Sahulland.
The \it{Lydekker Line} runs along the border of the Sahul continental shelf.
\it{Weber’s Line} is variously described as the area where the two types of fauna (from Sundaland
and Sahulland) are thought to be equally mixed,
or as the easternmost extent of certain categories of fauna from Sundaland (such as reptiles).
The flora (plants) do not show such clear patterns.
Wallacea is the transition zone between the two biogeographical regions (the zone between the Wallace
Line and the Lydekker Line),
and is a zone that requires multiple deep water crossings for animals or humans to traverse
\citep[58]{Bellwood2007,Bellwood2019ed,OConnor2015,Oppenheimer2003}.
We return to this in \crf{ch:MorLes}.

Wallacea includes flora and fauna from both Sundaland
and Sahulland, as well as flora
and fauna endemic (native to,
unique) in the region such as the \it{babirusa} (\it{Babyrousa babyrussa;} \frf{02-02})
and the \ili{Buru} red lory (\it{Eos bornea cyanonothus} \frf{02-RedLory}).
Wallacea is where Austronesian-speaking peoples coming from the north or west would have first encountered
marsupials, and developed a need to talk about them \citep{Blust1982,Blust2012b,Schapper2011}.
Biogeographical Wallacea includes Sulawesi,
Lombok, and western \ili{Sumbawa}.
“Linguistic Wallacea” does not include Sulawesi,
Lombok, or western \ili{Sumbawa}.
Linguistic Wallacea excludes these western parts
and extends further eastward to encompass parts of Indonesian Papua, including SHWNG languages.
Our study is focused on linguistic Wallacea minus SHWNG
(see Figures \ref{fig:01-05}, \ref{fig:01-06}, \ref{fig:02-01}).

%\begin{figure}
	%\includegraphics[width=0.4\columnwidth]{02-02.jpg}
	%\fcp{Babirusa (\it{Babyrousa babyrussa})}{02-02}
%\end{figure}

\begin{figure}\centering
  \subfigure[Babirusa (\it{Babyrousa babyrussa}) ]{\label{fig:02-02}
    \includegraphics[width=0.45\textwidth]{02-02.jpg}
    }\hspace{15mm}
    \subfigure[\ili{Buru} red lory (\it{Eos bornea cyanonothus})]{\label{fig:02-RedLory}
    %\subfigure[\ili{Buru} red lory]{\label{fig:02-RedLory}
    \frame{\includegraphics[width=0.25\textwidth]{02-RedLory.jpg}}
    }
  \caption{Fauna unique to Wallacea}
\end{figure}

Humans travelling from Asia to reach Australia
and New Guinea, regardless of whether they came from the north or west,
would have had to make multiple deep water crossings in Wallacea.
\citet[1]{Bellwood2019b} points out,
“No land bridges ever crossed the full expanse of this region of deep seas
and steeply plunging coastlines during the Pleistocene Ice Ages,
even during glacial maxima,
and humans migrating from Taiwan or Sundaland towards New Guinea
and Australia always had to cross sea gaps between islands to reach their goals.” 

At the time of initial human encroachment into the region,
\citet[215]{OConnor2015} is even more specific,
“the global sea level was 60--70 m lower than it is today.
While this may have reduced distances between islands,
at l\ili{east} eight separate crossings,
and one leg of at l\ili{east} 70 km,
would have had to be made,
regardless of the route taken” to reach Australia.
(See \frf{02-04}.)

\begin{figure}
	\includegraphics[width=\columnwidth]{02-04.jpg}
	\fcp{Deep water crossings in Wallacea}{02-04}
\end{figure}

\citet{BellwoodIrwin1998} find evidence of human habitation in north Maluku around 35,000 BP (BP = before present).\footnote{
		Technically BP is set at 1950,
		when the term “before present” came into use.}
At greater depths,
around the 40,000 BP range,
evidence of human habitation has not been found there \citep{Bellwood2019b}.

\citet[217]{OConnor2015} reports that in modern day Timor-Leste,
“at two sites, Jerimalai shelter and Lene Hara Cave,
[human] occupation dates back to 42,000 cal BP”. %(cal = calibrated).
She also reports:

\begin{quote}
	``More than thirty-nine thousand fish bones from twenty-three
	different taxa were identified overall from this test pit.
	The most common fish were Scombrids (tuna),
	Scaridae (parrotfish), Carangidae (trevallies),
	Balistidae (triggerfish), and groupers (Serranidae).
	Other major fish families identified are Lethrinidae (emperors),
	Lutjanidae (snappers), unicornfish (Acanthuridae),
	wrasses (Labridae), and puffers (Tetradontidae).
	Rays and sharks (Elasmobranchi) are also present \citep{OConnorClarkson2011}.
	Next to fish, marine turtles made up a large part of the diet.
	[{\ldots}] Although fish remains are most plentiful in Holocene levels,
	the earliest Pleistocene occupation level (42,000
	and 38,000 years ago) contained significant quantities of fish bone.
	[{\ldots}] Fast-swimming pelagic species [deep sea,
	open ocean], such as tunas, account for almost 50 percent
	of the fish bones in these levels.''
	\citep[219]{OConnor2015}
\end{quote}

She notes that to reach Australia or New Guinea from Sundaland,
“Whatever route was taken,
this was a remarkable accomplishment at this early date
and demonstrates the capacity for complex planning
and technological innovation that is the hallmark of our species.
First, it required the successful navigation of multiple water crossings between the stepping-stone islands en route.
Second, the colonists arriving on each new island had to adapt to unfamiliar fauna,
flora, and landscapes” \citep[215]{OConnor2015}.

\begin{figure}\centering
	\includegraphics[width=0.4\columnwidth]{02-03.jpg}
	\fcp{Skipjack tuna (\it{Katsuwonus pelamis})}{02-03}
\end{figure}

The implications of the large numbers of pelagic
fish bones (see \frf{02-03}) are highly significant for
our understanding of the prehistory of the region.

\begin{quote}
	``The pelagic fish bones from Jerimalai clearly demonstrate that the level of maritime capacity
	and fishing technology of the early modern humans
	who colonised our region was much higher than previously believed.
	[{\ldots}] Pelagic species are fast-moving,
	travel long distances, and are mostly caught offshore.
	They are therefore quite difficult to capture,
	as compared with inshore fish.
	[{\ldots}] Whatever method was used,
	the high proportion of pelagic species at Jerimalai indicates
	that their capture was probably planned rather than fortuitous.
	{\ldots} The earliest direct evidence of fishing technology in the East Timor sites is the tip of a hook
	made on Trochus sp. shell found in Jerimalai in a level dated to between 23,000
	and 16,000 years ago \citep{OConnorClarkson2011}.
	Indeed, this currently stands as the earliest evidence for the manufacture of fishhooks anywhere
	in the world.'' \citep[220f]{OConnor2015}
\end{quote}

She concludes:

\begin{quote}
	``The early modern humans moving out of mainland Asia must already have possessed a high degree of maritime
	capacity, which would have enabled them to make water crossings into Island South\ili{east} Asia in the first place.
	Settlement in the islands of Wallacea would have presented many new challenges,
	not l\ili{east} of which would have been the restricted range of terrestrial animal species on many of them.
	[{\ldots}] by 42,000 cal BP,
	the people living in this region were skilled in the acquisition of marine resources,
	which requires high levels of planning.
	They caught fast-moving pelagic fish such as tuna,
	as well as sharks,
	rays, and a range of inshore fish; to do this,
	they must have possessed a range of sophisticated organic technology.''
	\citep[221]{OConnor2015}
\end{quote}

We do not know what languages these earliest human travellers into
and through the region spoke.
We do know that they were not Austronesian languages at that early time.
We do know that early humans did reach
and settle the New Guinea region far to the \ili{east},
requiring many deep water crossings.
We do not know if they had a continuous presence throughout our target region or only in some places.
For example, \citet[220]{OConnor2015} says, “Jerimalai shelter appears to have little evidence of occupation
38,000--24,000 years ago {\ldots} [but] nearby Lene Hara Cave was occupied during this period.”\footnote{
		The relative gaps in dating are interesting,
		and probably reflect the scarcity of archeological digs in the Wallacea region.
		Presumably early humans travelling to Sahulland (the Australia
		and New Guinea region) would have had to traverse Wallacea.
		Dating human habitation in Sahulland currently ranges
		from 50,000--60,000 BP \citep{OConnellAllen2003,ReddStoneking1999}.
		Yet the earliest dates in Wallacea seem to be around 42,000 BP.
		One assumes that additional digs will find evidence of earlier dating,
		or possibly evidence of other routes to Sahulland.}

By the time the Austronesians arrived in the region
tens of thousands of years later, around 2,000--1,500
BC (\citeauthor{Bellwood2007} [\citeyear[123]{Bellwood2007}]; ≈ 4,000--3,500 BP),
it appears that the pre-existing populations already in place spoke various non-Austronesian languages
of at l\ili{east} two different language families that are loosely combined under the term ‘Papuan’.

Before returning to discussion of the pre-existing non-Austronesian languages of eastern
Indonesia and Timor-Leste, it is helpful to draw attention to a few points
regarding evidence of the arrival of Austronesian-speaking peoples.

\begin{itemize}
	\item There is abundant evidence of the introduction of various flora (particularly
				food crops) and fauna associated with the arrival of the Austronesian-speaking peoples
				into the region around 4,000-3,500 BP \citep{Bellwood2007}.
	\item There is evidence of “Neolithic settlement of northern Maluku”,
				associated with the Austronesian-speaking peoples dating about 3,500 BP \citep[514]{Spriggs2011}.
				More specifically “plain red-slipped pottery,
				polished stone adzes, pigs,
				and dogs dating from 3300 BP in the Uattamdi rockshelter”
				assumed to be associated with Austronesian-speaking
				peoples have been found on Kayoa Island \citep[218]{Bellwood2019a}.
	\item There are traces of “pig-bone [{\ldots}] red-slipped pottery,
				1 sherd with incised decoration; chert
				and obsidian flakes” associated with Austronesian-speaking peoples at a site in the \ili{Banda} Islands in
				central Maluku dating around 3,000 BP \citep[520--521]{Spriggs2011}.
	\item It is entirely possible that the Austronesian incursion(s) into northern Maluku are separate from
				and independent of the incursions into the Lesser Sundas,
				southern and central Maluku (see \srf{15-06}).
	\item There was a period of intense,
				rapid, and far-reaching migrations by the Austronesian-speaking peoples,
				covering thousands of islands
				and thousands of kilometres spanning just a few centuries.
				\citet[520--521]{Spriggs2011} says,
				“The spread of Neolithic \tsc{An}-speaking cultures across much of ISEA
				[Island South\ili{east} Asia] is a similar phenomenon,
				in terms of its rapidity,
				to the Lapita expansion beyond the Bismarck Archipelago between about 3100
				and 2900 BP when that culture spread beyond Near Oceania
				through the south-\ili{east} Solomons,
				Vanuatu, New Caledonia, Fiji
				and into western Polynesia.
				[{\ldots}] Polynesia, which was first occupied only
				about 3500 BP in the case of the Marianas
				and possibly Palau, and not long after about 3100 BP
				for the rest of Island Melanesia
				and western Polynesia.”
\end{itemize}

So, we note that Remote Oceania appears to have been settled by people speaking Austronesian languages
around a broadly similar time period as for the settlement of eastern Indonesia.
It is not clear if this rapid expansion into the Pacific involved or by-passed what is now central
and southern Maluku, Timor-Leste,
and the rest of Nusa Tenggara (the Lesser Sundas).
The assumption based on the distribution of EMP languages (including SHWNG) has been that the jump-off
point for what later became Oceanic
and Polynesian languages somehow involved north Maluku.
But that view may no longer have strong support from archaeology \citep[4]{Bellwood2019b}.
We return to this issue in \crf{ch:MorLes}.

We do know that many grammatical features commonly associated with Oceanic languages
\citep{Pawley1973,Ross1998} are also found in many languages
of present day eastern Indonesia and Timor-Leste.
These include the grammar and semantics of active
and non-active verbs, verbs of motion
and posture, serial verbs,
prepositional verbs, complex possessive constructions, alienable (indirect)
and inalienable (direct) possession.
And, for example, an edible possessive construction commonly associated with some Oceanic languages,
is also found in \ili{Selaru} in southern Maluku \citep[47--48]{Coward2005}.
Phonological features associated with Oceanic languages such as PMP *-ay > \it{e},
*-aw > \it{o}, and *ə > \it{o} are found sporadically in different parts of the target region.
But these are also found widely in languages in Sulawesi
and other languages west of the Blust line.

In trying to see where the linguistic picture
and the archaeological picture link up,
\citet[511]{Spriggs2011} comments: 

\begin{quote}
	``We can use the spread of the ISEA Neolithic as a proxy for AN language spread,
	as justified at length by \citet{Pawley2004}
	and \citet{Ross2008}, among others.
	\it{In doing this, it is very hard to see anything between PMP
	and EMP at all from the archaeology.} It would seem that movements out of Taiwan were rapid after about
	4000 BP and by 3800 BP dialects of PMP were spoken everywhere from the Philippines to eastern Borneo,
	Sulawesi and south to East Timor,
	spreading with the first pottery-using cultures in those areas.
	Currently the dates for the EMP area in northern Maluku do seem to reflect a later time of spread,
	at about 3500 BP as with Palau
	and the Marianas and Java.'' \citep[511]{Spriggs2011}	[emphasis ours]
\end{quote}

So, in terms of the limited archaeological data currently available,
Spriggs sees that the Austronesian incursion into the
Timor region occurred several hundred years before
the evidence of Austronesians in northern Maluku,
precisely during a time period where a span of
200--300 years appears to have been incredibly significant
for the dispersal of Austronesian-speaking peoples.\footnote{
		Given the scarcity of archaeological digs
		in the region, this may be making much ado about nothing.} 

There are questions about where
and when and how Austronesian-speaking peoples came
into the target region that are introduced briefly
here, and addressed in greater detail in \crf{ch:MorLes}.
\citet[98]{Blust2008} says,

\begin{quote}
		``\ili{As} noted in \citet{Blust1993},
		\ili{Central} Malayo-Polynesian almost certainly was the product of rapid movement from north\ili{east} Indonesia
		through the central and southern Moluccas and into the Lesser Sunda islands,
		forming a dialect chain that favored the ready diffusion of innovations across incipient language boundaries,
		and hence destroying any tree-like structure that may have previously arisen.'' \citep[98]{Blust2008}
\end{quote}

While he provides no evidence for such movement,
he nevertheless criticises \citet{Schapper2011} for an “implicit claim” that both the glosses in the
languages and the limited distribution of bandicoots show the Austronesian-speaking peoples encountered the marsupial
cuscus before they encountered the bandicoot,
saying her scenario “implies an entry via Sulawesi into the Lesser Sundas,
with subsequent eastward movement into the Moluccas” \citep[274]{Blust2012b}.
While that goes beyond what she actually says,
it is nevertheless revealing.
Since Blust rejects an incursion via Sulawesi,\footnote{
		We show in \crf{ch:Tal} that there is evidence
		from the historical phonology for at l\ili{east} one incursion into Maluku from Sulawesi.
		\citet[130]{Bloyd2020} also observes,
		“Sula was a casual swim’s distance from Sulawesi during the height of the Ice Age.”
		In chapters \ref{ch:Mar}
		and \ref{ch:MorLes} we also show several lexical ties between south\ili{east} Sulawesi,
		Maluku, and the rest of Wallacea.}
		we have to assume by “north\ili{east} Indonesia” (\citeyear[98]{Blust2008})
		he means Halmahera or Indonesian Papua.
Yet the limited archaeological evidence currently available has the Austronesian-speaking peoples reaching
Timor (presumably from the west),
300-500 years before they show up in Halmahera.
The scattered and disjointed pockets of conservative features in the historical phonology,
along with common direction of sound changes presents a much more complex picture than that proposed by Blust.

The scattered and uneven pockets of complex morphology along with non-cognate morphology from subgroup
to subgroup that does not trace back to PMP,
interspersed with other pockets of fairly simple morphology,
also paints a much more complex picture than that proposed by Blust.
And in a context of great mobility along with intense
and prolonged contact with typologically different languages,
the principle of “the area with the greatest diversity probably represents the homeland” cannot be assumed to hold true.
So while Blust’s scenario of Austronesian incursion into our target region does not seem to be grounded
in archaeology, historical phonology, morphosyntax, or the lexicon, we find that the linguistic evidence favours multiple
incursions, perhaps from different starting points,
and does not necessarily rule out his proposed scenario as one of several possibilities.
Yet we are also not actually sure what data or evidence his proposed scenario is based on.
We favour a both-and approach over an either-or approach on Austronesian incursions into the region.
We return to this issue in \crf{ch:MorLes}.

\citet{Spriggs2011} likewise comments on the broad patterns of contact by the Austronesian-speaking peoples
as they encountered other populations,
and implications for methodology.

\begin{quote}
	``Roger Green’s \citeyear{Green1991,Green2000} model of intrusion,
	integration and innovation captures the situation well in ISEA as well as the western Pacific,
	whether we are talking of material culture, language or people.
	There was indeed some migration out of Taiwan \citep{Kayser2008};
	there was mass recruitment of people from populations already resident in ISEA
	and Near Oceania as the Neolithic spread \citep{Soares2011};
	artefacts and practices were integrated from already-resident groups
	and others were discarded by them; new ideas were brought into being as unexpected human
	and environmental situations were encountered.'' \citep[523]{Spriggs2011}

	``[The] Neolithisation of ISEA was a new process of identity formation that seized the imagination of a
	mass of people on hundreds of islands across thousands of kilometres of ocean,
	spreading like a pulse across ISEA
	and into the Pacific over a few centuries.
	It spread through processes both of migration and recruitment in-place.''
	 \citep[524]{Spriggs2011}

	``[Consequently:] we are struggling to find appropriate models to deal with just what happens during such
	temporal ‘hot spots’.'' \citep[523]{Spriggs2011}									
\end{quote}

Contact complicates comparative-historical linguistics.
And contact is not a trivial add-on,
but rather a core feature of our target region,
and one that is necessary to understand for unravelling
the complexities of the data in eastern Indonesia and Timor-Leste.
Many people are finding the traditional branching tree model (see \frf{01-02})
and an assumption of a monolithic spread of the
Austronesian-speaking peoples to be inadequate to account
for the complexities of the region.
\citet[9]{Klamer2019} summarises the issues:

\begin{quote}
	``Indeed, for the micro-groups that have been studied reasonably well (cf.
	\citealp{Adelaar2005}), the absence of evidence for intermediate
	ancestral nodes connecting them suggests a complex dispersal scenario.
	There may have been just a few MP groups migrating into ISEA,
	but then each of these groups did not spend a sufficiently long period of time together as unified speech
	communities after their split from MP.
	Because if they had,
	they would have innovated words in their lexicon or developed different pronunciations which would have
	been inherited in their daughter languages
	and could now be used to define the subgroup as such by its unique lexical and phonetic innovations.
	There may also have been not a few but \it{many} different MP communities migrating at different times
	and in different directions throughout history.
	Overall, the lack of shared inherited features suggests that most of the MP languages that dispersed
	into ISEA shared a relatively short local micro-history.
	In contrast, the shared inherited features of,
	for example, proto-Oceanic, indicate that this group did remain a unified speech community for a period
	that was sufficiently long to create such defining innovations
	and pass them on to their daughter languages.
	In sum, there is no evidence that the dispersal of MP speakers in ISEA followed simple trajectories;
	instead, it is more likely that multidirectional processes of contact
	and language admixture have been at work for millennia.''
	\citep[9]{Klamer2019} [emphasis in the original]
\end{quote}

Returning now to the Papuan languages,
we note that those still spoken in the region today include the \tsc{Timor-\ili{Alor}-Pantar} (TAP) languages
(\citealp{Klamer2017,HoltonRobinson2017} -- see \frf{02-05}),
and the \tsc{West Papuan} languages of North Halmahera
(\citealp{Voorhoeve1988,GrimesGrimes1994,HoltonKlamer2018} -- see \frf{02-06}).
But the implications of \citet{Donohue2007} are also that we have evidence of earlier Papuan languages
spoken as far west as \ili{Sumbawa} into the early 1800s.
In this case, the \ili{Tambora} language was wiped out,
not by the Austronesian-speaking peoples,
but by an erupting volcano.
There is widespread trace evidence of various sorts of contact with SOV (Papuan) languages found
in Austronesian languages that are hundreds of kilometres from any Papuan languages spoken today,
such as on Flores,
Rote, \ili{Buru}, Seram, \ili{Kei}, Tanimbar, and \ili{Babar}.
We elaborate on these issues in \srf{02-03-01}.

\begin{figure}[ht]
	\includegraphics[width=\columnwidth]{TimorPapuan.pdf}
	\fcp{Austronesian and Papuan languages in greater Timor}{02-05}
\end{figure}

\begin{figure}[p]
	\centering
	\includegraphics[width=0.87\columnwidth]{HalmaheraPapuan.pdf}
	\fcp{Austronesian and Papuan languages in north Maluku}{02-06}
\end{figure}

There are interesting questions about the prehistory
of the non-Austronesian languages in the region.
(\frf{02-05} and \frf{02-06}.)
There are several obvious possibilities:

\it{Scenario {\#}1:} The non-Austronesian languages are remnants of the ancestral languages spoken by
the earliest humans on the way from Island South\ili{east} Asia to the island of New Guinea.
These ancestral languages were perhaps the linguistic precursors of the many Papuan language families
that later diversified significantly.

\it{Scenario {\#}2:} The non-Austronesian languages today are the result of more recent westward moving
back migration(s) out of the New Guinea mainland.
The assumption in this scenario,
based partly on the spread from the New Guinea region of several important food crops
and animals, is that they were already in place in the Timor region before the arrival of the Austronesian-speaking
peoples around 4,000--3,500 BP.
This view is reflected in \citet[85]{Blust2013} with regards to the TAP languages,
“All non-AN languages of
the Lesser Sundas have been assigned to the Trans-New Guinea Phylum of Papuan languages,
and evidently represent a prehistoric back-migration from the western end of New Guinea to the islands
of \ili{Pantar}, \ili{Alor} and adjacent parts of Timor.”\footnote{
		In spite of accepting such a scenario of a prehistoric back-migration
		from western New Guinea,
		\citet[40]{Blust2009} appears sceptical that “the reference to
		‘contact-induced change’ draws on several publications
		by Donohue that advance the novel theory that a Papuan-based
		trading network spanned the \ili{Malay} archipelago
		prior to the arrival of AN speakers during the second millennium BC.”}
There is a growing body of typological
and lexical evidence of a connection between TAP
and \tsc{West Bomberai} languages which are thought to link to \tsc{Trans New Guinea} languages
\citep{Windschuttel2018,UsherSchapper2018,UsherSchapper2022}.

\it{Scenario {\#}3:} a combination of scenarios {\#}1 + {\#}2.
While there is increasing evidence in support of scenario {\#}2,
that nevertheless does not negate the possibility of scenario {\#}1.
That is, the more recent westward back-migration from the New Guinea region intermixed with a remnant
Papuan substrate from the early human migrants who left evidence in the archaeology.

The reality is actually much more complicated than either {\#}1 or {\#}2.
Given the span of time of over 40,000 years,
and the evidence of early boat technology that supported deep water crossings
and open ocean fishing,
it is unrealistic to limit non-Austronesian incursions to one time,
in one place and from one direction.

For example, there is evidence of at l\ili{east} two different waves of non-AN languages in \ili{Alor}-\ili{Pantar}.
Several groups claim their ancestors came from the Timor mainland.
The Papuan language enclave of \ili{Oirata} on \ili{Kisar} Island \ili{east} of Timor also traces their origins to a relatively
recent migration from the Timor mainland.
Some of these could have happened at the
same time or after Austronesian contact
(\citealp[8f]{GrimesTherikGrimesJacob1997}, \citealp{Klamer2017}).

But there are many other complexities to add into the mix as well.
Scenario {\#}1 above is problematic,
in that the break-up of the \ili{Alor}-\ili{Pantar} (AP) branch of TAP appears to have occurred about the same time
as the arrival of the Austronesian-speaking peoples.
\citet[10]{Klamer2017} says,

\begin{quote}
“the level of lexical
and grammatical similarity in the AP family does not support
an age of more than several millennia,
and the reconstructed vocabulary of proto-AP appears
to contain Austronesian loan words such as ‘betel
nut’ \citep{HoltonSchapper2012,Robinson2015}.
Ancient Austronesian loans found across the \ili{Alor}-\ili{Pantar} family
following regular sound changes suggest
that the AP family split up after being in contact
with the Austronesian languages in the area.”
\end{quote}

Presumably
the higher-level relationships of the TAP
family predated the arrival of the Austronesians.
In \srf{02-05} we discuss potential methodological
problems with linking the timing of the break-up
of the AP branch of TAP to the arrival
of the Austronesian-speaking peoples.

An issue that supports scenario {\#}2 above is growing evidence of westward movement out of the New Guinea
region of certain food crops whose botanical origins come from there,
and that the dispersal was happening prior to the arrival of the Austronesian-speaking peoples
\citep{DenhamDonohue2009,DonohueDenham2010,Schapper2017}.
For example, “the initial stages of banana
and sugarcane domestication with subsequent diffusion westward [out of New Guinea] occurred before the
appearance of Malayo-Polynesian languages in ISEA” \citep[20]{DonohueDenham2010}.
They are specifically referring here to banana cultivars developed from \it{Musa acuminata} spp.
\it{banksii}, and to \it{Saccharum officinarum} (sugarcane),\footnote{
		Since sugarcane (\it{Saccharum
		officinarum}) is now known to have originated from the New Guinea region,
		it is assumed that the \tsc{PAn} reconstruction of *CəbuS ‘sugarcane’ either reflects a very early spread north
		to Taiwan, or a different \it{Saccharum} species \citep[515]{Spriggs2011}.
		Very likely this northern variety was \it{Saccharum sinense}.}
but their discussion also includes taro, breadfruit and sago.
This dispersal of food crops westward from the New Guinea region has huge implications for methodology
that we return to later.

With regard to genetic DNA research in the region,
\citet[18]{Klamer2019} summarises recent findings:

\begin{quote}
	``Genetic research confirms that during the Holocene period in ISEA,
	groups with Asian ancestry spread through ISEA.
	Although the data on ISEA are still very sparse,
	there is presently statistical support that Taiwan played
	an instrumental role in Holocene dispersals
	across ISEA, but the data do not clearly support a \it{majority}
	genetic input from there.
	[{\ldots}] \citet{HudjashovCox2017} show
	that during the period in question (<4,000 years)
	there is very limited evidence for autosomal
	gene transfer to Indonesian populations,
	and that which is tends to be from the Philippines.
	Moreover, there is no evidence of gene transfer to western Indonesia,
	and Taiwan does not seem to have a direct role.
	Overall, the speakers of Austronesian languages today
	are very diverse in biological terms,
	and the genetic evidence indicates that the evolutionary processes
	in ISEA have been very complex
	and have involved a great deal of admixture.
	Finally, the genetic data are often inconclusive about the directions
	of these dispersals'' \citep[265]{LansingSudoyo2011}.
	[emphasis in the original]
\end{quote}

\citet[22]{Fricke2019} contributes additional information
relevant to the region:

\begin{quote}
	``There appears to be a clear boundary between Asian genomes in the west
	and non-Asian genomes in the \ili{east} between the islands of Bali
	and Flores, more specifically between \ili{Sumbawa}
	and Flores 
	(\citealp[1837f]{KarafetHammer2010}, \citealp[2447]{HudjashovCox2017}).
	This genome boundary coincides with the Wallace line,
	a boundary of transition between Asian
	and Australian types of flora and fauna.
	[{\ldots}] The sampling results from the eastern part of Flores (Bama),
	from Lembata and from Timor show an almost half-half division of the Papuan
	and the Austronesian genome component.
	This genetic profile is very similar to the sampling results on the islands of \ili{Alor}
	and \ili{Pantar} which show only a slightly higher Papuan component.
	Towards the west, the central
	and western part of Flores (Bena
	and Rampasasa) and Sumba show a smaller Papuan component than found further \ili{east}.
	On Sumba, the Austronesian component is the biggest.
	On Flores, the admixture is more diverse.'' \citep[22]{Fricke2019}
\end{quote}

It is important to remember that genes
and languages are not inherently linked to each other.
For example, in urban Mexico there are hundreds of
thousands of people of genetic Uto-Aztecan or Mayan
ancestry who speak Indo-European Spanish
and not their heritage languages.
Languages move across genetic populations.
Both individuals and whole societies can shift
to speak languages unrelated to their DNA ancestry.
So the DNA evidence noted above is useful for confirming
that different gene pools have mixed in the
target region, and that there is a genetic boundary
that roughly coincides with the Blust line,
but is not otherwise very useful linguistically.
To highlight a few key points:

\begin{itemize}
	\item There is evidence of human habitation in our target region dating back to 42,000--35,000 BP.
				These humans had boats or other technology that allowed them
				to make multiple deep-water crossings between
				islands that were tens of kilometres distant,
				and to catch large numbers of open ocean (pelagic) fish species.
				We do not know what languages these early inhabitants spoke.
				Given the time depth
				and the technology and the yearly \ili{east}
				and west monsoons, there is no reason to assume there was only one journey by one group at one time in one direction.
	\item There were groups speaking a variety of Papuan languages scattered across our target region
				at l\ili{east} as far west as eastern \ili{Sumbawa},
				prior to the arrival of the Austronesian-speaking peoples in the region around 4,000--3,500 BP.
	\item Also prior to the arrival of the Austronesian-speaking peoples,
				several important food crops were dispersing \it{westward} out of the New Guinea region.
				These included a species of banana,
				sugarcane, taro, breadfruit and sago.
				The spread of terms such as {\#}muku ‘banana’ across unrelated language families raises serious questions
				for the region about subgrouping arguments that are based on the scattered distribution of certain lexical items.
	\item When the Austronesian-speaking peoples arrived in the region,
				again there is no reason to assume there was only one journey by one group at one time in one direction.
				They had the technology to go in any direction at any time,
				seas and weather permitting.
				The current limited archaeological evidence has a 300 -- 500 year time difference between evidence for
				Austronesian-speaking peoples in the Timor region (earlier),
				and in North Maluku (later),
				which is quite different from the scenario proposed in \citet[98]{Blust1993,Blust2008}.
				\ili{As} we show in this study,
				the complexity and discontinuities in the linguistic data
				in the region suggests either very rapid incursion(s),
				or the possibility of multiple incursions,
				possibly at multiple times,
				possibly from several directions,
				from different starting points,
				and different regions where they got off their boats.
	\item We infer that the persistence over three or four millennia of Papuan languages in the region to
				the present day, along with structural, typological
				and lexical accommodation in both directions between Austronesian
				and Papuan languages, suggests a contact scenario of bilingualism
				and intermarriage over centuries that also resulted in some language shift,
				rather than necessarily a dominance model of the “farming/language dispersal” hypothesis.
				Warfare seems to have been mostly localised over land,
				women, and power, rather than large-scale territory-grabbing efforts.\footnote{
						See also \citet{Schapper2020a}
						that large-scale warfare in eastern Indonesia
						and Timor-Leste seems to post-date European contact.} 
	\item The human gene pool of our target region is “very complex
				and [{\ldots}] involved a great deal of admixture”,
				with a genetic transition near the Wallace line
				(\citealp{Oppenheimer2003}, \citealp[22]{Fricke2019}).
				While this is important for understanding the region,
				it does not tell us about the languages or their prehistory.

\end{itemize}

\scp{Evidence of on-going movement and contact}{02-02}

The information in this section is not intended to be exhaustive,
but merely a sample to illustrate some of the complexities that
need to be understood in the region.
The information is broadly organised by topics,
and within each subsection there is an attempt to be geographical
and chronological where possible.\footnote{
		A reviewer suggested that this section be reorganised chronologically.
		However, there are too many unknowns
		and disagreements among the experts about when many of these
		issues involving contact in the region began,
		to be able to do so coherently.
		The earliest spice trade had a different sociolinguistic
		impact on the region than the much later European quests -- but
		they both involve the spice trade,
		so there is some necessary redundancy in our discussion.
		While there is clear evidence that trade in cloves spans at l\ili{east}
		two millennia, there is some evidence pushing it back to almost four millennia.
		Sandalwood trade may span two millennia or more.
		There could have been some overlap with the Lapita trade.
		How does one sequence those unknowns chronologically?
		It is not known how far back bird-of-paradise trade goes.
		It is not known when the tripang (sea cucumber) trade began,
		and there is also disagreement.
		Technologically it could have been going on for centuries before
		the 300 year-old threshold touted by some.
		The presence of loanwords such as ‘money’
		and ‘chain’ from Austronesian languages into Australian Aboriginal
		languages does not allow us to identufy when the trade began,
		but is merely evidence that the contact continued into
		the historical present.}

\sscp{Trade networks that span centuries and even millennia}{02-02-01}
There has been much trade through the region,
it has involved people from outside the region,
and it has been happening for a long time -- some of it spanning two millennia or more.

The Lapita cultural complex associated with Oceanic languages in the Pacific involved
trade spanning thousands of kilometres
and several centuries, perhaps beginning around 3470--3250 cal.
BP \citep[95]{SpechtTerrell2014}.
The Lapita trade involving pottery
and obsidian in near Oceania also had links back to Borneo
and to the northern Philippines -- and was thought to have links to Timor and Maluku.
\citet{SpechtTerrell2014} explain these Lapita connections as follows:

\begin{quote}
	``Interisland movement of goods also has a long history within ISEA,
	as obsidian in late Pleistocene
	and mid-Holocene contexts on Timor,
	Borneo, and the \ili{Talaud} Islands almost certainly derived from off-island sources
	\citep{AmbroseReepmeyer2009,Reepmeyer2011,Reepmeyer2011}.
	Contacts with the New Guinea region in the early
	and mid-Holocene are attested by the recovery of marsupial faunal remains in prepottery
	contexts in Timor and the \ili{Talaud} and northern Molucca Islands
	\citep{FlanneryNitihaminoto1995,FlanneryBalasubramaniam1998}.

	The first clear evidence for contact between ISEA
	and the [Bismarck] Archipelago comes from the transfer of pottery technology,
	followed by the commensal animals discussed above.
	In the reverse direction,
	New Britain obsidian occurs at Bukit Tengkorak on Borneo in contexts contemporary
	with Middle Lapita
	\citep{BellwoodKoon1989,Bird1989}.''
	\citep[116--117]{SpechtTerrell2014}

	``The Borneo connection with the [Bismarck] Archipelago continued after Lapita.
	Two undated double-spouted pots found in the \ili{Manus} Islands show ‘‘a singular similarity’’
	\citep[p. 28]{Kennedy1982} to vessels of uncertain age in Niah caves in Sarawak \citep{Harrison1971},
	and a piece of bronze of probable ISEA origin found on
	\ili{Manus} is assigned a post-Lapita date \citep{Ambrose1988}.
	In summary, there is evidence in pre-Lapita times of long-term interaction between
	ISEA and New Guinea and between New Guinea
	and the Archipelago, but not for direct interaction between ISEA and the Archipelago.''
	\citep[117]{SpechtTerrell2014} 
\end{quote}


It is important to understand that these are pointing to \it{trade} links involving
travel in two directions spanning a long period of time.
This provides early direct evidence of things like obsidian associated with the Lapita
cultural complex, sourced in the Bismarck Archipelago north
and \ili{east} of New Guinea,
having been brought westward to Borneo,
going against the flow of the singular idea of a monolithic
eastward migration by Austronesian-speaking peoples.
And this trade involved Sulawesi and Borneo.
There appears to have been travel back-and-forth,
and contact maintained over a significant period of time.
The trade may not have involved end-to-end trade routes,
but rather worked through intermediate links -- much like the spread
of {\#}muku ‘banana’ (see \srf{02-03-01-07}).

Earlier ideas that these Lapita links somehow involved north Maluku are being questioned by some.
\citet[219]{Bellwood2019a} states:

\begin{quote}
	``Actual cloves have been found in early to middle first millennium AD contexts at
	Segaran IIA in West Java
	and Mantai in Sri Lanka \citep{ManguinIndradjaja2011}.
	Furthermore, Indian pottery in quite large quantities is now well known from the
	sites of Sembiran and Pacung in northern Bali,
	dating to circa 2000 years ago or just before
	\citep{ArdikaBellwood1991,ArdikaYuliathi1997,Calo2015}.
	Neither ancient cloves nor Indian pottery have been excavated in the Moluccas,
	but it seems relatively certain that Indian contact could \it{potentially} have reached
	the Moluccas by 2000 years ago.
	So could, and did,
	the technological developments that are widely associated by archaeologists,
	linguists, and historians with a combination of Indian,
	Mainland South\ili{east} Asian, early Malayo-Chamic,
	and Chinese trading acumen.
	These include metal items
	and metallurgy, Dong Son drums, glass beads,
	nephrite ornaments (although these are not yet found in the Moluccas),
	and probably also some new Y-chromosome haplogroups \citep{Kusuma2016}.''
	\citep[219]{Bellwood2019a} [emphasis in the original]
\end{quote}

The later so-called “spice orgy” (e.g. \sqcitet{Wright2007}{Wright2007} among others)
associated with Europe
and eastern Indonesia has impacted the history of the world,
including Columbus bumping into America on his way to find a back door to the Spice Islands of Maluku.
It also has a role in the histories of the Netherlands,
Ceylon, Capetown, Manhattan, Suriname,
and economies of islands such as Granada in the Carribean,
and Zanzibar in Africa.
(See \sqcitet{Milton1999}{Milton1999} as one of many sources.)

Cloves (\it{Syzygium aromaticum} [formerly \it{Eugenia caryophyllatta}]),
nutmeg and mace (\it{Myristica fragrans}) are native to northern
and central Maluku respectively (see \frf{02-07} to \frf{02-10}).
Sandalwood (\it{Santalum album Linn}.) is thought to be indigenous to Timor
and Sumba and the Indian subcontinent.
Trade in these items goes back hundreds,
and even thousands of years,
long before European presence in the region.

\begin{figure}
	\centering
	\frame{\includegraphics[width=0.3\columnwidth]{02-07.jpg}}
	\fcp{Drying cloves for shipping}{02-07}
\end{figure}

Although the exact trade routes are not known,
\citet[67]{Lapian1965} reports finds of cloves in present day Syria in the remains of a
pantry room in a common household,
dating to around 1,700 BCE.\footnote{
		While the botanical identity of the Terqa cloves
		has been questioned, the issue has not been settled,
		nor does it detract from other evidence of clove trade spanning over 2,000 years
		(nearly 4,000 years if the Terqa cloves are indeed cloves).}
\citet[1--2]{Andaya1993} notes
that some of the earliest written references to cloves include a Chinese source during
the Han dynasty (206 BC-222 AD).
Pliny the Elder describes cloves reaching the Roman world by the first century AD.
\citet[67]{Lapian1965} finds clear evidence that cloves had reached Europe by the 4th century,
when Silvester, the Bishop of Rome from 314--335,
received a gift of 150 pounds of cloves.
\citet[78]{Lapian1965} also notes that before Europeans gained direct access to the source
of the spices in Maluku in 1512,
that Asian traders brought the spices from Maluku to the western part of the archipelago
where they were then purchased by Arab,
Chinese and other foreign traders
and eventually made their way to China,
the Middle East, and Europe via Venice.
There is some evidence that many of these “Asian” traders were from Java
and \ili{Makasar} and were associated with the coming of trade \ili{Malay} into the eastern regions,
already functioning as a language of wider communication by the time of the arrival
of the Europeans \citep{Grimes1991,Grimes1994}.

%%%%%%%%%%%%%%%%%%%%%%%%%%%%%%%%%%%%%%%%%%%%%%%%%%%%%%%%%%%%%%%%%%%%%%%&%
%%                                                                     %%
%% NOTE FOR TYPESETTING: I have added some manual spacing to have      %%
%% the pictures not too close together, but also not too far apart     %%
%% and to have them appear on the same page as the text referencing    %%
%% them. The sizes have also been reduced (they were too big in        %%
%% the proofreading version.                                           %%
%%                                                                     %%
%% Please consult before editing how these pictures have been typeset! %%
%% Changing the sizes could have bad flow on effects for pagination.   %%
%%                                                                     %%
%%%%%%%%%%%%%%%%%%%%%%%%%%%%%%%%%%%%%%%%%%%%%%%%%%%%%%%%%%%%%%%%%%%%%%%%%

\begin{figure}\centering
  \subfigure[Fresh picked cloves \it{(Syzygium aromaticum)}]{\label{fig:02-08}
    \frame{\includegraphics[width=0.31\textwidth]{02-08.jpg}}}\hspace{2mm}
  \subfigure[Preparing nutmeg and mace for drying]{\label{fig:02-09}
    \frame{\includegraphics[width=0.31\textwidth]{02-09.jpg}}}\hspace{2mm}
  \subfigure[Drying nutmeg \it{(Myristica fragrans)}]{\label{fig:02-10}
    \frame{\includegraphics[width=0.31\textwidth]{02-10.jpg}}}
\fcp{Spices native to eastern Indonesia}{02-08--02-10}
\end{figure}

%\begin{figure}
  %\subfigure[Bronze Dong-Song drum from Vietnam used as bridewealth in \ili{Alor}]{\label{fig:02-11}
		%\includegraphics[width=0.25\textwidth]{02-11.jpg}}\hspace{10mm}
  %\subfigure[One of many  varieties of tripang \it{(Holothuroidea sp.)}]{\label{fig:02-12}
		%\includegraphics[width=0.25\textwidth]{02-12.jpg}}\hspace{10mm}
	%\subfigure[Sign at ``Macassan   Beach'' in north\ili{east}   Arnhem Land, northern Australia]{ \label{fig:02-13}
		%\includegraphics[width=0.25\textwidth]{02-13.jpg}}
%\caption{Long-distance trade\label{fig:02-11--02-13}}
%\end{figure}


%\begin{figure}
    %\subfigure[Fresh picked cloves  \it{(Syzygium aromaticum)}]{\label{fig:02-08}
		%\frame{\includegraphics[width=0.25\textwidth]{02-08.jpg}}
%}
    %\subfigure[Preparing nutmeg  and mace for drying]{\label{fig:02-09}
		%\frame{\includegraphics[width=0.25\textwidth]{02-09.jpg}}
    %}
    %\subfigure[Drying nutmeg   \it{(Myristica fragrans)}]{\label{fig:02-10}
		%\frame{\includegraphics[width=0.25\columnwidth]{02-10.jpg}}
%}
	%\fcp{Spices native to eastern Indonesia}{02-08--02-10}
%\end{figure}

Timor and Sumba figured in the sandalwood trade.
\citet[1]{Hagerdal2012} gives the first historical
mention of the sandalwood trade on Timor
from around 1350 \tsc{ce} in Chinese sources.
Since contact with the Indian subcontinent dates back to the
first and second centuries \tsc{ce},
it is possible that sandalwood trade in Timor
and Sumba could pre-date the first mention in Chinese sources
by several centuries.

There were early trade networks linking Vietnam
and Thailand to the islands of eastern Indonesia.
\citet[17]{Klamer2019} says, 

\begin{quote}
``During the Metal Age ({\tl}3,000--1,500 BP),
traces of trade networks become more visible.
For example, Galipaud (in press) mentions three types of links: one linking south
Vietnam and Thailand and the Indonesian islands,
the other linking northwest Vietnam,
Laos and Yunnan to Sumatra,
eastern Indonesia and New Guinea, and the third linking south Vietnam with the Philippines and Serawak.
South Vietnam and Thailand are linked with the Indonesian islands way into the \ili{east},
to the Lesser Sunda islands Flores,
Lembata, \ili{Pantar}, and \ili{Alor} by the existence of Dong-Song drums
and the existence of cemeteries with burials where bodies were put in certain types of jars.
The transport of some Dong-Song drums into ISEA is one of the clearest traces of
a regional connection between the SE Asia mainland
and the islands and is dated after the dispersal of MP [Malayo-Polynesian] languages.
The cemeteries with urn containers are always coastal
and therefore suggestive of mobile communities connected by a vast network of social
relations extending possibly as far as the South\ili{east} Asian mainland.
In a recent excavation of the cemetery of Pain Haka in eastern Flores
\citep{GalipaudBuckley2016}, the skeletons were dated
and revealed that urn burials were already in use by 2,800 BP in Pain Haka,
and that the cemetery remained in use for more than 500 years,
without any evidence of imported metal.
Pain Haka excavations provide clear evidence of burial practices similar to those
documented in other parts of SE Asia.
Chief among these is the use of pottery jars alongside other forms of container for
the interment of the dead.
The dating of the site combined with the fact that this burial practice is present
over such a wide geographic area suggests a widespread belief system during the Neolithic
period across much of SE Asia.
In other words, burial traditions linking ISEA with the South\ili{east} Asian mainland
started well before the Metal Age,
i.e., before the dispersal of MP languages.'' \citep[17]{Klamer2019}
\end{quote}

Tripang trade (sea slug,
sea cucumber; \it{Holothuroidea spp}.; see \frf{02-12})
and \it{Trochus} (top shells) used in making buttons
and jewellery also brought traders from many other parts of Asia into our target region.
They were also a catalyst for early contact with northern Australia (particularly
Arnhem Land) by seafarers from what is now central (including Makassar)
and southeastern Indonesia (including the Lesser Sundas and Maluku).
Scores of boats made the yearly trek to northern Australia,
planted tamarind trees, and left archaeological sites,
pottery, words, and part of their gene pool.
The term “Macassan” in northern Australia is a grab bag to refer to anyone from the
Indonesian archipelago, and is not limited to Makassar.
Hundreds of loanwords from Austronesian languages have been identified in Aboriginal
languages of northern Australia \citep{WalkerZorc1981,Evans1992}.
Grimes has been told stories by Aboriginal people in both \ili{east}
and west Arnhem Land of their grandfathers
and great-grandfathers accompanying boats on the return journeys northward
and marrying local women in central
and eastern Indonesia, some returning to their homelands after decades away.

Yolŋu Matha words like orthographic ‹\it{dhimurru}› [ˈd̪imuru] ‘\ili{east} wind,
\ili{east} monsoon’ and ‹kapala› [ˈkapala] ‘ship’ with initial stress reflect the
telltale extrametrical echo vowels on roots ending with \it{r,
l, s} that narrow the probable source languages to Makassar,
Konjo, or \ili{Selayar} in South Sulawesi \citep[96]{GrimesGrimes1987}.
Words like ‹djalataŋ› [ˈʤalataŋ] ‘south’
and ‹rranhdhiŋ› [ˈran̪d̪iŋ] ‘chain’ are probably from some form of \ili{Malay} (cf.
\ili{Malay} \it{səlatan} ‘south’; \it{rantai} ‘chain’).
But ‹bärra› [ˈbara] {\tl} [ˈbaːra] ‘west wind,
west monsoon’ and ‹djarra› [ˈʤara] ‘horse’ could be from many possible sources;
for example many languages in the Lesser Sundas have \it{jara} [ˈʤara] ‘horse’,
whereas the final nasal of Makassar \it{ʤarraŋ} ‘horse’ is commonly preserved in
Yolŋu Matha languages and so is unlikely to be the source.

%%%%%%%%%%%%%%%%%%%%%%%%%%%%%%%%%%%%%%%%%%%%%%%%%%%%%%%%%%%%%%%%%%%%%%%%
%%
%% NOTE FOR TYPESETTING: I have added some manual spacing to have
%% the pictures not too close together, but also not too far apart
%% and to have them appear on the same page as the text referencing
%% them. The sizes have also been reduced (they were too big in
%% the proofreading version.
%%
%% Please consult before editing how these pictures have been typeset!
%% Changing things could have bad flow on effects for pagination.
%%
%%%%%%%%%%%%%%%%%%%%%%%%%%%%%%%%%%%%%%%%%%%%%%%%%%%%%%%%%%%%%%%%%%%%%%%%

\begin{figure}\centering
  \subfigure[Bronze Dong-Song drum from Vietnam used as bridewealth in \ili{Alor}]{\label{fig:02-11}
		\includegraphics[width=0.25\textwidth]{02-11.jpg}}\hspace{10mm}
  \subfigure[One of many  varieties of tripang \it{(Holothuroidea sp.)}]{\label{fig:02-12}
		\includegraphics[width=0.25\textwidth]{02-12.jpg}}\hspace{10mm}
	\subfigure[Sign at ``Macassan   Beach'' in north\ili{east} Arnhem Land, Australia]{ \label{fig:02-13}
		\frame{\includegraphics[width=0.25\textwidth]{02-13.jpg}}}
\caption{Long-distance trade\label{fig:02-11--02-13}}
\end{figure}

%\begin{figure}
    %\subfigure[Bronze Dong-Song drum from Vietnam used as bridewealth in \ili{Alor}]{\label{fig:02-11}
	    %\frame{\includegraphics[width=0.3\textwidth]{02-11.jpg}}
    %}
  %\subfigure[One of many  varieties of tripang \it{(Holothuroidea sp.)}]{\label{fig:02-12}
	%\frame{\includegraphics[width=0.3\textwidth]{02-12.jpg}}
	%}
    %\subfigure[Sign at ``Macassan   Beach'' in north\ili{east}   Arnhem Land, northern Australia]{ \label{fig:02-13}
	%\frame{\includegraphics[width=0.3\columnwidth]{02-13.jpg}}
	%}
%\caption{Long-distance trade\label{fig:02-11--02-13}}
%\end{figure}

In a recent article,
\citet{Schapper2022} has identified a handful loanwords in indigenous languages of
northern Australia which originate from languages within our target region: terms
related to ‘hawksbill turtle’ from reflexes of Proto-\ili{Aru} *kaloba,
terms meaning ‘fishing line’ from \ili{Aru} languages,
‘sea cucumber’ from words related to Proto-Rote-\ili{Meto} *ɓanafi,
and ‘tobacco’ from words related to Proto-Rote-\ili{Meto} *maɗo \citep[9--12]{Schapper2022}.
Schapper posits that such forms indicate that peoples from Wallacea may have undertaken
their own voyages to northern Australia.

%\largerpage
The long-standing seasonal interaction between the Austronesian-speaking peoples
and the people of Arnhem Land is still remembered
and highly valued, even though the Australian government has attempted to cut off
the relationship since Federation in 1901.
For example, on 14 June 2020 the Aboriginal band East Journey performed a concert
singing their songs in the \ili{Gumatj} language at the Bawaka homeland,
%\footnote{
		%See \url{https://www.facebook.com/VisitEastArnhemLand/videos/3812185228851950/?_rdc=1&_rdr}
		%for a recording of the concert.}
as part of a series of Covid-19 online concerts
from several bands in East Arnhemland.
Around 6:58 minutes into the concert they were singing a traditional song about
\it{liya-wayma kamboŋ malogo, liya-wayma kamboŋ dhua},
which in \ili{Malay} would be something like
\it{layar-layar kampuŋ Maluku, layar-layar kampuŋ tua}
‘the sails of Maluku village, the sails of the old village’.
In a contemporary song at around 7:54 minutes into the concert they sing
\it{ga wiripunydja kamboŋ dhua} “and another way/name is \it{kampuŋ tua} [old village]”,
with the next line clarifying “these speak of Bawaka”
(the remote homeland where the concert was filmed).
Several involved with the concert series have names
like Djakapurra (Singapore) and Djawa (Java).
There are songs about places like \it{Gupaŋa} (\ili{Kupang}).\footnote{
		Our appreciation
		to Concert Producer and Sound Engineer Andrew Grimes for narrowing down the phrases for us.
		He gave additional examples,
		but what is presented here suffices to make the point.}
On 9 July 2020 a Facebook
post from Numbulwar, further south in the gulf of Carpentaria,
was seeking assistance for translating a number of names associated with flags
and dances linked to the “Macassans”.
The memory and the significance of the contact is not lost.

While hundreds of loanwords have been identified going from Austronesian languages
into various Australian Aboriginal languages,
nobody seems to have seriously investigated loans going in the other direction.
But until we have better documentation of those languages,
the task would remain difficult.

Coastal west Papua and the \ili{Aru} islands supported a
long-standing trade in bird-of-paradise
specimens and feathers, and other exotic fauna.
This also brought a steady stream of traders from the western parts of the Indonesian
archipelago and further afield \citep{Andaya2017}.

In addition to the well described spice trade routes to \ili{Ternate},
\ili{Ambon} and \ili{Banda}, there was also a trade network from
\ili{Banda} to \ili{Onin} on the west coast of Papua.
The language of \ili{Geser}-\ili{Gorom} has also been used as a regional lingua franca around
eastern Seram and the west coast of Papua for generations
\citep{LoskiLoski1989,WalkerWalker1991,Zobel2014}.
These regional connections are also reflected in clan names on \ili{Ambon},
such as Latupapua ‘king of Papua’ (pers. comm. Zeth Latupapua).
These factors must be kept in mind when unravelling the linguistic
affiliations of \tsc{Seram-Tanimbar-Bomberai} languages.
There are telltale differences in the sound correspondences.
(See \crf{ch:STB}.) 

\citet{Rutter1930} describes seasonal slave raiders travelling from the southern
Philippines through northern Maluku to the Bird’s Head area of Papua,
before continuing on to Timor
and Rote through central Maluku,
and eventually returning north to the Philippines. Traces are widespread.
\citet[74]{Grimes2000a} describes stories on \ili{Buru} about
the \it{geba lano} ‘\it{lano} people’
slave raiders, and the local \ili{Buru} \it{epkitan}
‘fighting champions’ who were to watch for them from
high spots along the coast
and fight to protect their own people when the slave raiders came.
The \it{geba lano} are variously thought in different parts of \ili{Buru} to be from Papua,
from Seram, or from \ili{Tobelo}
and \ili{Galela} (Halmahera, north Maluku; West Papuan languages).
The term \it{lano} may be associated with the \ili{Ternate} (West Papuan language) term
\it{kolano} which refers to a political position
in societies of north Maluku \citep{Andaya1990},
or it may refer to Seram which is often referred to in central Maluku languages as
\it{selano} (as in the clan names Pattiselano or Huliselan, etc; [pers. comm.
Jules Pattiselano; pers. comm. Mus Huliselan]).
But most likely it refers to the seasonal Ilanun
slave raiders from the southern Philippines \citep{Rutter1930}.
Their impact is still felt
and feared after generations.

\begin{figure}
	\centering
	\frame{\includegraphics[width=0.5\columnwidth]{02-14.jpg}}
	\fcp{Traders from south\ili{east} Sulawesi entering \ili{Ambon} harbour}{02-14}
\end{figure}

While the role of the ships
and sailors from \ili{Makasar}
and Java is better known in the historical literature,
inter-island trade in eastern Indonesia into the present has,
for generations, been carried out by sea-going traders variously
associated with the islands of Buton,
Binongko, \ili{Wanci}, \ili{Cia-Cia}, \ili{Tukang Besi}
and others off south\ili{east} Sulawesi (see \frf{02-14}).
There are long-standing settlements from south\ili{east} Sulawesi scattered
throughout the region, including northern coastal Flores,
\ili{Pantar}, Obi, \ili{Buru} and Seram.
For example, \citet{Grimes2000a} notes that many communities along
the western and southern coasts of \ili{Buru} are long-standing settlements totalling
around 12,600 people from south\ili{east} Sulawesi.
Much of the knowledge
and terminology associated with the sea in the \ili{Buru} language,
such as ways to treat a stonefish puncture,
is said to have been learned from these settlers.
There is some evidence from the lexicon
and phonology (such as initial
and medial prenasalised stops) that indicates the recently extinct
language of \ili{Hukumina} in northwest \ili{Buru} was possibly a satellite
community of speakers from Sulawesi \citep{Grimes2010b}.

All of these trade links serve to illustrate the kinds of “contaminants”
that can muddy the waters in trying to unravel the phonological
prehistory of the Austronesian languages in the region.

\sscp{Empire outliers}{02-02-02}
Indian influence via the Majapahit Hindu-Buddhist empire (around
1293--1517) in \ili{east} Java left emissaries in \ili{Ternate}, \ili{Ambon}
(e.g. \it{Rehatta} clan; pers. comm.
Sonja Rehatta) and Timor (e.g.
the \it{Kase} clans; pers. comm.
Simon Kase); a caste system along the southern islands as far
\ili{east} as \ili{Kei}; and kin group names in places such as \ili{Kei}
and \ili{Buru} independently associated with the Majapahit empire
(e.g. \it{Masbait} clan, pers. comm.
Raja Anton Lesnussa). Javanese words are still used as the names
of islands in the region (e.g.
\it{Wetan} ‘\ili{east}’ and \it{Alor} ‘north’), and in names for introduced
fauna (e.g. \ili{Buru} \it{menʤaŋan} ‘deer’).

While there was some trade associated with the Majapahit presence
in eastern Indonesia, it is the empire,
its infrastructure and dominance over vassal states that has
left a lasting legacy in the region.
While Majapahit was the last of the Indic
kingdoms in Indonesia, Indian influence across
Indonesia traces back as early as the
first to second centuries CE.
See also \citet{Hoogervorst2016,Hoogervorst2017,Hoogervorst2023,Hoogervorst2024}
for the influence of South Asian languages in the region.

The people of Sabu often express their perception that their
identity and culture are linked to India.
\citet[199--244]{Fox1977} explores many parallels in the livelihoods
and economies of Rote
and Sabu with southern India
and Sri Lanka, all based around the lontar palm (\it{Borassus flabellifer}).

The political power-house of Gowa (Makassar,
in south Sulawesi, west of the Blust line
and not a “CEMP” language) dominated the Flores
and \ili{Sumbawa} region up until 1727,
leaving linguistic traces of varying degrees \citep[585]{VerheijenGrimes1995}.

After 1727, \ili{Bima} (a “CEMP” language of eastern \ili{Sumbawa}) dominated
the nearby islands around Flores.
This contact may account for some of the similarities behind
the undefined notion of a “\ili{Bima}-Sumba” subgroup,
for which \citet{Blust2008} has not been able to find support
in the historical phonologies.

\citet{vanFraassen1987} and
\citet[15--31, 146, 159--160]{Grimes1993}
describe the history of contact
between \ili{Ternate} and central Maluku (Sula,
\ili{Buru}, \ili{Manipa}, \ili{Kelang}, Boano,
Seram), and the influence from Ternatan ‘governors’ on the political
structures, social titles and language on \ili{Buru}.
The names of several individual Ternatan emissaries are known.
\ili{Ternate} is a Papuan language from north Halmahera (see \frf{02-06}).
\ili{Buru} is an Austronesian language in central Maluku.
The “memory” of the period of Ternatan control on \ili{Buru},
stopped by the Dutch in the mid-1600s,
is still talked about on \ili{Buru} today
and figures prominently in oral histories of the island.
Several of the Ternatan legacy positions/titles are still used
by some kin groups in the mountains.

While a number of \ili{Ternate} loans (a Papuan language) have been
identified in the \ili{Buru} language,
such as \it{matgugul} ‘title of \ili{Ternate} governor’,
\it{soa} ‘clan’ and \it{jou} ‘lord’, the evidence is much more
striking in the Sula Islands north of \ili{Buru}.
Fortgens was the investigator for the Holle lists for \ili{Soboyo},
\ili{Taliabu}, and \ili{Kadai} \citep{Stokhof1980b}.
In the \ili{Soboyo} data Fortgens identified 52 loans from \ili{Ternate}
and 4 from Java.
In the \ili{Kadai} data he identified 43 loans from \ili{Ternate}
and 3 from Java.
\citet{Collins1982c} has multiple references of \ili{Ternate} influence
and presence throughout Seram
and the islands in the \ili{Manipa} Straits,
as well as noting the presence of several \ili{Tobelo}
and \ili{Galela} speaking villages (both Papuan languages from north
Halmahera) around coastal western and northern Seram.
\ili{Ambon} \ili{Malay} \it{kusu} ‘cuscus’ is very likely a loan from a \tsc{North
Halmahera} language, such as \ili{Ternate} \it{kuso} ‘cuscus’,
particularly since no other language in central Maluku has a
form remotely similar in form (see \crf{ch:Mar}, particularly \srf{14-03-02}).

\sscp{Colonial greed and forced resettlements}{02-02-03}
The events described below
and their impact on language have been written about in more detail by
\citet{Grimes1994,Collins2003b,MusgraveEwing2006,Grimes2010b}.

Nutmeg (\it{Myristica fragrans}) is native to the \ili{Banda} Islands.
In 1621, Dutch \it{Vereenigde Oost-Indische Compagnie} (VOC)
forces decimated perhaps 90{\%} of the native population in order
to gain control of the islands
and the trade in nutmeg.
Remnants of the \ili{Banda}-speaking population fled to the south\ili{east}
to the \ili{Kei} islands,
and set up the refugee communities of \ili{Banda}-Elat
and \ili{Banda}-Eli that remain today,
albeit with a very different pattern of language contact than they had formerly.

The Hoamoal wars of 1651--1656 in central Maluku were an effort
by the Dutch VOC to gain a monopoly in the spice trade,
and limit planting to keep prices high.
The VOC took leaders hostage from many coastal communities,
and forcibly resettled whole or partial language communities
around their forts, creating multiethnic, multilingual communities around several
points of concentration (e.g.
\ili{Batumerah}, \ili{Manipa}, \ili{Kayeli}; see \frf{02-15}).
\citet[17]{vanFraassen1983} describes it as follows:

\begin{quote}
		``The war ignited by the 1651 rebellion [against the Dutch] continued until 1656.
		The rebels received Macassarese support,
		and the war was in no way restricted to Hoamoal [west Seram].
		Arnold de Vlamingh van Oudtshoorn,
		the Dutch commander-in-chief in this war,
		was also launching attacks on the VOC’s opponents elsewhere,
		among other places in \ili{Buru} and \ili{east} Seram.
		[{\ldots}] The Ambonese region was formally removed from Ternatan
		control and the institution of the Ternatan governorship was abolished.
		Hoamoal, [{\ldots}] was completely depopulated [{\ldots}] The
		chiefs of Hoamoal had a place of residence assigned to them in
		\ili{Batumerah}, in the vicinity of the VOC’s chief fortress in \ili{Ambon}.
		The population of the islands of Boano,
		\ili{Kelang} and \ili{Ambelau} was transferred to the island of \ili{Manipa},
		where the Company had a small fort.
		All the Muslim chiefs of the coastal areas of \ili{Buru} were obliged
		to settle in the neighbourhood of the Company’s fort at \ili{Kayeli}.
		[{\ldots}] The evacuated areas were systematically destroyed
		and rendered unfit for reoccupation.
		Contracts were concluded with the chiefs of \ili{Buru}
		and North and East Seram in which the latter conceded themselves
		to be subordinate to the VOC; promised to entertain no relations
		with other nations or rulers,
		to keep out all Macassarese,
		Malays, and Javanese.'' \citep[17]{vanFraassen1983}
\end{quote}

\begin{figure}
	\centering
	\frame{\includegraphics[width=0.45\columnwidth]{02-15.jpg}}
	\fcp{Entrance to VOC fort at Kayeli}{02-15}
\end{figure}

Note the repeated references to alliances with “Macassarese, Malays, and Javanese”.
The fallout from these events continues to shape the sociolinguistic
landscape of central Maluku today,
nearly 400 years later.
The last speakers of several of these transplanted language communities
have just died out in the last three or four decades,
and many of their descendants have completed the shift to other
languages such as \ili{Ambon} \ili{Malay}.

\sscp{Internal migrations}{02-02-04}
Oral histories throughout the region indicate continuous
cross-fertilisation with regional migrations both pre- and post-European contact.
For example, local origin myths tie \ili{Ambelau} (south of \ili{Buru}) together
with \ili{Nusalaut} (south of Seram).
There are clans in several language groups around Timor that
claim their origins trace to Seram,\footnote{
Seram written with
a final m{\#} is from a Portuguese orthographic tradition that
writes final nasals as ‹m›.
Locally in central Maluku it is pronounced [ˈseran] or more frequently [ˈseraŋ].
\ili{Gorom} is similarly pronounced [ˈgoroŋ] around central Maluku.
On Timor, orthographic ‹Tetum› is pronounced [ˈtetʊn].
Compare the pronunciation of Portuguese orthographic \it{Belém}
[beˈleŋ], \it{sim} ‘yes’, etc., which are also not bilabial.}
such as the \it{Seran} clans in the \ili{Tetun}-speaking region (pers. comm. E. Seran;
pers. comm. Father Dr. M.
Bria Seran), and another in the \ili{Helong}-speaking region.
It remains to be seen if these are merely folk etymologies or
if there is more to them.
Given the complex history of the region,
it is not beyond possibility.

In most regions of eastern Indonesia there are origin groups
(kin groups, clans) whose origin stories are local.
They believe they are autochthonous,
originating from the island where they now are.
The origin stories of other groups have a second type of people
originating elsewhere and migrating in the distant past to their current locations.
Yet these latter second-wave groups are also fully assimilated
into the local social structures
and languages, and may even be given the role of a leader or
king \citep{Fox1995,Grimes2006}.

For example, \citet[77--84]{Grimes1993} describes a number of \it{autochthonous}
origin groups on \ili{Buru} whose stories have their founding ancestor
or ancestress appearing out of a rock,
a tree or a stream on the island.
The closer to the headwaters,
the more senior in precedence in relation to other origin groups
(but since each group governs itself,
precedence does not give them any governing rights over another
group; it only provides a script for how elder-younger social relationships should interact).
The further downstream the origin story,
the more junior in precedence.
The place where the founding ancestor built the first house is
known as the \it{tean elen} ‘place of planting the house pole into the ground’.

Other \it{non-autochthonous} origin groups came from elsewhere,
settled on \ili{Buru}, learned the language,
intermarried, and became assimilated into \ili{Buru} culture.
The location where their founding ancestor or ancestress first
set foot on \ili{Buru} is known from their origin stories,
and is called the \it{waga eʔnohon elen} ‘place of disembarkation
from the boat’ (from /en-toho-n/ with regular participial metathesis
from \it{toho} ‘descend, disembark’ < PMP *tuRun ‘descend’),
or an \it{esnaba} ‘disembarkation’.
For example, the \it{Mual} kin group attributes its origins to
the arrival of their ancestress,
Bokisraja, from the Hoamoal peninsula of Seram (note the phonological
similarity between \it{Mual} and the final part of Hoa\emph{moal}).
Their \it{esnaba} is near the village of Masnana.
The Nalbessy kin group \it{esnaba} is near the village of Wae Hotong.
Neither of those groups live near their \it{esnaba}.
Other kin groups claim to have come from other unnamed islands.
But power is tied to place,
so in 1990 when a non-autochthonous kin group that had migrated
to lake \it{Rana} (< PMP *danaw ‘lake’) in the centre of the
island had a ‘bad death (many people die in a short span of time)’,
they abandoned their village at \ili{Rana} where they had lived for
many generations and migrated four days walk to the northwest coast of \ili{Buru} to
reconnect with the power associated with their ancestral \it{waga
eʔnohon elen} near Bara to try
and restore harmony and prosperity.

\citet{Grimes1993} also discusses how these \it{non-autochthonous}
kin groups are categorised as \it{geb fuk Buru} ‘people of \ili{Buru}
island’, and not as \it{geba man lau tau-n} ‘outsiders’ (lit: people whose
origins flow from the sea),
such as the migrants from Sula
and Sulawesi that have been on \ili{Buru} for 12--15+ generations.
Similar patterns are found throughout the region.

The entire north coast of \ili{Buru} has been settled for 15+ generations
by around 12,000--13,000 people from the Sula islands to the
north, completely dominating speakers
of the \ili{Lisela} language on \ili{Buru} \citep{Grimes2000a}.
This has been largely ignored in attempts to establish a linguistic
link between \tsc{Buru} and \tsc{Sula}, which share only about 27{\%}
of their vocabulary
(even less between \tsc{Buru}
and \tsc{Taliabu}), and have very different historical phonologies,
although they do share PMP *R > \it{h}.
(See chapters \ref{ch:Tal} and \ref{ch:SB}.) 
The west coast of \ili{Buru} has been settled by migrants from south\ili{east}
Sulawesi since recorded history.
Their long-standing dominance has resulted in loans into the
\ili{Buru} language, and in things like the corruption of place names such as \ili{Buru}
\it{Kak Pala Madat} (the name of the highest mountain in central
and southern Maluku at 2,428m; \it{kak < kaku} ‘mountain,
ridge’ also found in other place names such as \it{Kak Tuan,
Kak Hanga}, etc.) to \ili{Sanskrit}-\ili{Malay} \it{Kapala Madan},
which recently became the formal government name for the new
regency in western \ili{Buru} where there are very few native \ili{Buru}
people to stand up for their heritage or tell their histories.
Several village names on coastal western \ili{Buru} have been similarly
corrupted or have outside names.
This has not happened nearly so much in areas where these immigrants
are not so dominant.

Volcanic activity around 1976 near the southern Maluku islands
of \ili{Teun}, \ili{Nila}, and \ili{Serua} (\tsc{Timor-Babar} languages)
triggered a government-sponsored
forced resettlement of the speakers of those three languages
to central Seram, hundreds of kilometres to the north,
which has severely impacted the vitality of these languages \citep{Edwards2025}.

What is described above is only a sample for illustrative purposes,
not an exhaustive list.
Some historical linguistic studies in eastern Indonesia have
overlooked these events that have triggered contact-induced language
change in the region -- including from other Austronesian languages.
Such contact issues must be sorted out before comparative work
based on word lists in isolation can be substantiated.



\scp{Complexity in the languages}{02-03}
In addition to normal language diversification resulting from
Austronesian dispersal, there were also widespread changes
that happened in the typology,
structures, and semantics due to contact with various Papuan
languages already in place (\srf{02-03-01}).
Numerous sociolinguistic factors also combined to result in a
plethora of synonyms in the region (\srf{02-03-02}).
Lexical borrowing can be fairly trivial.
Structural and typological changes from contact are quite significant.
	\sscp{Evidence of contact-induced change}{02-03-01}
\ili{As} shown in the following chapters,
under the surface it appears that historically there were multiple
incursions by Austronesian-speaking populations into the region,
who came into contact with pre-existing populations speaking
non-Austronesian (Papuan) languages that were typologically predominantly SOV.
There was bilingualism and intermarriage over millennia,
which resulted in Papuan languages influencing not only the vocabulary
(which by itself can result from relatively superficial contact),
but also the structures
and semantics of Austronesian languages, and \it{vice versa}.
Having both island and mountain ecologies, both tropical rainforest
and semi-arid climates, and newly encountered flora
and fauna in the region all meant that patterns of contact
and diffusion are quite different from island group to island
group, regardless of whether the contact was with Austronesian or Papuan speaking populations.
For example different patterns of localised contact provide the
most likely explanation for \ili{Havu}
and \ili{Dhao} being relatively closely related in their historical
phonologies and parts of their lexicon,
yet very different in their syntax and typologies \citep{Grimes2010a}.
The local contact histories in northern \ili{Buru},
northern Seram and northern Flores are significantly different
from those in southern \ili{Buru},
southern Seram and southern Flores respectively
\citep[11]{Grimes2000a,GrimesTherikGrimesJacob1997}.

\citet{ThomasonKaufman1988,Schapper2015,Klamer2019}, and
\citet{GrimesCIP} point out that the kinds of contact-induced
structural and typological changes found in languages of the region
and discussed in this section are best explained as reflecting
intense contact over a long period of time.
\citet[14--16]{Klamer2019} walks carefully through the logic of how
linguistic features that are not found in Austronesian languages
in the west and north, but are commonly found in both the Papuan languages
and the Austronesian languages of our target region can be attributed
to contact with the Papuan languages.
It is particularly important to understand the significance of
Austronesian languages that are typologically SVO shifting away
from structures that are commonly associated this typology,
and taking on structures that are commonly associated
with SOV typologies, but sometimes are only rarely associated with SVO typologies.
This is highly significant for understanding the linguistic prehistory of the region.

There are a variety of typological and lexical features in the languages
\ili{east} of the Blust line that are not typically found in languages to
its west and its north. These features tend to have a scattered
and/or non-uniform distribution meaning they cannot be attributed
to a single proto-language, or often even to multiple proto-languages.
Instead, these features appear to have been acquired by contact
after the break-up of proto-languages.

So for instance, the preposed possessor introduced in \srf{01-05}
is the normal pattern throughout our target region.
Yet within the \tsc{\ili{Bima}-Lembata} subgroup both preposed
and postposed possessors are found, with the preposed possessor
the most common order within the \tsc{Flores-Lembata} node,
yet even within \tsc{Flores-Lembata} post possessors also occur.

Similarly, in the lexicon \citet{DenhamDonohue2009,DonohueDenham2009,Schapper2015}
have pointed out that the pseudo-etymon {\#}muku `banana'
has a Papuan source and is frequent in the southern region
of our target region. But it co-occurs with reflexes retaining PMP *punti `banana'
(as well as other diverse etyma) even among members within the same node of a subgroup.
So, for instance, the \tsc{Timor-Babar} language \ili{Tetun} has
*punti > \it{hudi} while its sister \ili{Habun} within the \tsc{Eastern Timor}
node has {\#}muku > \it{muku}. Similarly, in the \tsc{Southwest Maluku}
node of \tsc{Timor-Babar} reflexes of *punti predominate,
but \ili{Kisar} has \it{muʔu}, while reflexes of **biak occur in
three other nodes of \tsc{Southwest Maluku}.
See \srf{02-03-01-07}, \frf{15-11}, and appendix \ref{ch:BanTer}
for more details.

Attempting to explain such distributions as the result of inheritance
from a common ancestor is futile. Instead, different subgroups, nodes,
and languages have each undergone contact in different ways and to different
extents which have resulted in the kinds of uneven distribution that we see.
A selection of the kinds of features that are found in such an uneven manner
is given in \trf{02-Features} divided into typological and lexical features.
This list is not intended to be exhaustive.
Some of these features are due to contact with Papuan languages,
while others may be due to contact between Austronesian
languages of the region after their arrival, and some may be due to both.
We discuss these features in more detail in chapter \ref{ch:MorLes}.

\begin{table}
	\centering\tcp{Contact features Austronesian languages of Wallacea}{02-Features}
	\st{6pt}\begin{tabular}{l@{ }ll}\lsptoprule
 \mc{2}{l}{Feature} & see \\	\midrule
 \mc{2}{l}{Preposed possessor} & \srf{02-03-01-01} \\
 \mc{2}{l}{Alienable-inalienable distinction} & \srf{02-03-01-01} \\
 \mc{2}{l}{Clause-final negation} & \srf{02-03-01-02} \\
 \mc{2}{l}{Postpositions} & \srf{02-03-01-03} \\
 \mc{2}{l}{Complex and non-decimal numerals} & \srf{02-03-01-04} \\
 \mc{2}{l}{absence of velar nasal /ŋ/} & \srf{13-01-01} \\
 \mc{2}{l}{subject prefixes on verbs} & \srf{13-03-03} \\
 \mc{2}{l}{Human/non-human distinction} & \srf{07-04-03} \\	\midrule
 *kəRaŋ & `shell, turtle' & \srf{13-01-05-03} \\
 **mantəR & `cuscus' & chapter \ref{ch:Mar} \\
 {\#}muku & `banana' & \srf{02-03-01-07} \\
 `head leaf' & = `hair' & \srf{02-03-01-05} \\
 `inside' & = `seat of emotions, character' & \srf{02-03-01-06} \\
 *duRi `thorn' &> *duRi `bone' & \srf{15-04-02-03} \\
 \mc{2}{l}{kinship systems and terminologies} & \srf{07-02-03}, \srf{15-04-02-04}, \srf{15-04-03-02} \\
 **malip & `laugh, smile' & \srf{15-04-04-01} \\
 **maya & `tongue' & \srf{15-04-04-02} \\
		\lspbottomrule
	\end{tabular}
\end{table}

The presence of a cluster of the features in \trf{02-Features}
in sets of languages adds to the complexity of unravelling what
is due to retention, shared innovation, parallel development,
and/or what is due to contact.

Although our discussion of contact in this work focusses on the influence of Papuan
languages on Austronesian languages in the region,
it has also been shown that Austronesian languages influenced
Papuan languages \citep{KlamerReesinkVanStaden2008,ReesinkDunn2018}.
For example, the development of an inclusive-exclusive distinction
in the pronouns of some Papuan languages is attributed to Austronesian contact.
And some see the shift to SVO typology in some nodes of Papuan
language families that have SOV typology in other nodes as due
to contact with Austronesian languages \citep[190--193]{Voorhoeve1988}.
The contact reflected by such changes is not trivial,
and imply far more significant contact than is required for
simply borrowing a few words.

	\sscp{The prevalence of synonyms}{02-03-02}
On a technical note before getting further into the discussion
of synonyms, some lexicographers are sceptical that there are very many exact
synonyms (true synonyms, full synonyms),
but accept that there may be many near synonyms
in a language (see \citealp{Zgusta1971}).
We are using the term ‘synonym’ here in a broad sense to refer
to multiple words or phrases for talking about the same or similar things.
Regardless of whether they are exact synonyms,
near synonyms, or only closely related meanings,
the nature of lexical taboos
and parallelisms means that many words that have slightly different
meaning in isolation are treated “as if” they are synonyms when
they are linked together for stylistic and sociolinguistic reasons.

In a development with many parallels to the discussion of PMP
*punti and {\#}muku (both meaning ‘banana’; see \srf{02-03-01-07}),
research summarised in \citet[260]{Schapper2011} notes that a species
of the marsupial cuscus (\it{Phalanger orientalis}) was introduced
by humans westward out from New Guinea into parts of Maluku
and Timor, although the Sulawesi ‘dwarf cuscus’ (\it{Strigocuscus
celenbis}) and the larger ‘bear cuscus’ (\it{Ailurops ursinus}) further
west on Sulawesi are apparently not from human transfer.
All are marsupials of the family \it{Phalangeridae}.
Like the food crops mentioned previously,
this happened several thousand years before the arrival of the
Austronesian speaking peoples into the region.

\citet{Blust1982} reconstructed two terms to
PCEMP, *kandoRa ‘cuscus’ and *mans(aə)r ‘bandicoot’, for marsupials
\ili{east} of the Wallace line.
\citet{Schapper2011} points out that the distribution of glosses
from most languages in our target region actually better supports
PCEMP *mans(aə)r meaning ‘cuscus’,
later only shifting to ‘bandicoot’ in Oceanic languages.
And the evidence supporting the full form of *kandoRa ‘cuscus’
outside of EMP is extremely limited in distribution to only one
language right near the geographical boundary of EMP,
and problematic for several other reasons (see \crf{ch:Mar}).
Reflexes of *kandoRa ‘cuscus’ are found in only a single node
of one of our Wallacean subgroups.

\citet{Blust2012b} acknowledges that Schapper is essentially
correct about the distribution
and glosses of *mans(aə)r as ‘cuscus’ outside of EMP.\footnote{
One of the remaining disputed glosses is discussed further in \srf{02-04-04}.}
He also concedes that “It is,
indeed, puzzling that reflexes of *kandoRa are not more widely
distributed in eastern Indonesia” (\citeyear[267]{Blust2012b}).\footnote{
		We also concede that the apparently conservative form of \ili{Watubela} \it{kadola}
		‘marsupial’ (just west of Indonesian west Papua) is puzzling
		and needs an explanation,
		particularly given that other conservative forms are found far
		away on the eastern end of New Guinea.
		We note that the trisyllabic \ili{Watubela} form is quite different
		in its vowels and structure from formally close words in \tsc{Seram-Tanimbar-Bomberai} and SHWNG languages.
		The comment in \citet{BlustTrussel2020} under the entry for *mansar ‘bandicoot’,
		dismissing \citeauthor{Schapper2011}’s (\citeyear{Schapper2011}) evidence that most languages outside
		of EMP have the gloss of ‘cuscus’ for cognates of *mansar,
		by saying her claim is “based largely on non-specific glosses
		in languages of eastern Indonesia,” lacks integrity.
		We address these issues more fully in \crf{ch:Mar}.}
But he is also dismissive of suggestions to revise his original (\citeyear{Blust1982})
reconstructed glosses and claims of the historical development.
He acknowledges that the two marsupial terms are “a critical
piece of evidence [{\ldots}] for the phylogenetic unity of CEMP”,
although not the sole evidence (\citeyear[262]{Blust2012b}).
The basis of some of his objections reveals his assumptions about
language and the Austronesian dispersal.

From Blust’s perspective, if we accept Schapper’s approach,
we are forced “to accept PCEMP *kandoRa
and *mans(aə)r as synonyms” (\citeyear[267]{Blust2012b}).
Yet, he says, “somewhat awkwardly,
this inference forces her to recognise both *kandoRa
and *mans(aə)r as competitors for the
meaning ‘cuscus’ in PEMP” (\citeyear[267]{Blust2012b}).
And, “we must ask why a language
would have synonyms for ‘cuscus’” (\citeyear[268]{Blust2012b}).

First, his comment is revealing in assuming “a [single] language”.
But the response to his question about synonyms is multi-faceted
and is well documented in the region on all counts.
They include: evidence of many other words with synonyms in the
target region, taboos, parallelisms and doublets, and different patterns of
contact in different parts of the same language.
But first we note that the reconstructed Austronesian vocabulary
in \citet{BlustTrussel2020} has many doublets,
synonyms, and near synonyms at all levels -- as
anything reflecting natural languages should.


		\ssscp{Many competing synonyms in the region}{02-03-02-01}

Many words have competing synonyms,
near synonyms or variants in the region.
For example, \citet[249]{Blust1993} observes that PMP *ka-labaw ‘rat’
and PCEMP *kanzupay ‘rat’ coexist \ili{east} of the Blust line.
\citet{BlustTrussel2020} list PMP *tiduR/*tuduR,
and PMP *qinəp/PCEMP *qenəp,\footnote{
		\citet{BlustTrussel2020}
		give PCEMP *qenəp as a doublet for PMP *qinəp ‘lie down to sleep’.
		However all their supporting forms are for SHWNG
		and \tsc{Oceanic} languages, suggesting that this form is more appropriate
		for PEMP, rather than PCEMP.
		While they list \ili{Kowiai} as a SHWNG language,
		we reclassify it as \tsc{Seram-Tanimbar-Bomberai} (see \crf{ch:STB}).
		That means there is only one reflex of *qenəp found outside of
		EMP languages in their data.}
in which the local glosses in
the many languages represented consistently hover around ‘sleep, lie down’.
\ili{As} an example of these terms in real languages,
\ili{Buru} has \it{bage} ‘sleep,
lie down’, \it{toro-n} ‘sleep soundly’ (from PMP *tuduR),
while the closely related \ili{Lisela} language on \ili{Buru} has \it{ine}
‘sleep, lie down’ (from PMP *qinəp).
There are additional forms listed in Blust
and Trussel for \tsc{PAn}
and PMP under ‘sleep’ as well.
It is common for there to be many competing ways to talk about the same thing.
The distribution of these competing terms (and many others) are
in languages in close proximity to each other.
(See the distribution maps throughout this work.)

\ili{As} another of many examples,
we have already noted that ‘banana’ has multiple generic terms
scattered through the region.
Reflexes of PMP *punti coexist alongside {\#}muku from Papuan
sources (\srf{02-03-01-07}).
\citet{DonohueDenham2009} document additional reflexes also extant
in the region, that all mean ‘banana’.
These are also found outside of our target region,
and must all be considered as competing synonyms for ‘banana’.
Data in \trf{02-06} are gleaned from \citet{DonohueDenham2009}; with our subgroups
and some corrections added.
The reconstructed forms are also modified in some cases
to better account for the reflexes observed in the region.

\begin{table}
	\centering\tcp{Reflexes of competing terms for banana}{02-06}
\begin{threeparttable}
	\begin{tabular}{lll}\lsptoprule
Language	&	\hp{*}‘banana’	&	Subgroup		\\ \midrule
	&	\hp{*}*kalauR	&				\\ \midrule
\ili{Bima}	&	\hp{**}kalo	&	\tsc{\ili{Bima}-Lembata}			\\
\ili{Komodo}	&	\hp{**}kalo	&	\tsc{\ili{Bima}-Lembata}				\\
\ili{Kambera}	&	\hp{**}kaluu	&	\tsc{\ili{Bima}-Lembata}				\\
\ili{Wanukaka}	&	\hp{**}kaloʔa	&	\tsc{\ili{Bima}-Lembata}				\\
\ili{Barakai}	&	\hp{**}kalaur	&	\tsc{Aru}			\\
SW \ili{Tarangan}	&	\hp{**}kalor	&	\tsc{Aru}				\\
	&	**kula	&					\\
\ili{Asilulu}	&	\hp{**}kula	&	\tsc{\ili{Ambon}-Seram}			\\
\ili{Laha}	&	\hp{**}kula	&	\tsc{\ili{Ambon}-Seram}				\\
\ili{Haruku}	&	\hp{**}kura	&	\tsc{\ili{Ambon}-Seram}				\\
\ili{Saparua}	&	\hp{**}kula	&	\tsc{\ili{Ambon}-Seram}				\\
\ili{Sou Amana Teru}	&	\hp{**}kula	&	\tsc{\ili{Ambon}-Seram}				\\ \midrule
	&	*telewa	&			\mc{1}{r}{also in SHWNG}		\\ \midrule
\ili{Nusalaut}	&	\hp{*}telewa	&	\tsc{\ili{Ambon}-Seram}				\\
\ili{Teluti} (\ili{Tehoru})	&	\hp{*}telewa	&	\tsc{\ili{Ambon}-Seram}				\\
East \ili{Tarangan}	&	\hp{*}taragwar	&	\tsc{Aru}				\\ \midrule
	&	*biak	&			\mc{1}{r}{also in SHWNG and Papuan lgs.}	\\ \midrule
\ili{Teun}	&	\hp{*}fiwa	&	\tsc{Timor-Babar}				\\
\ili{Nila}	&	\hp{*}hia	&	\tsc{Timor-Babar}				\\
\ili{Ambelau}	&	\hp{*}bia-e	&	\tsc{Sula-Buru}				\\
\ili{Lisela}	&	\hp{*}fia-t\su{†}	&	\tsc{Sula-Buru}				\\
\ili{Sanana}	&	\hp{*}fia	&	\tsc{Sula-Buru}				\\
\ili{Mangole}	&	\hp{*}fia	&	\tsc{Sula-Buru}				\\
\ili{Kadai}\su{‡} &	\hp{*}fia	&	\tsc{Taliabu} (\tsc{Celebic}) 			\\
\ili{Taliabu}	&	\hp{*}fia	&	\tsc{Taliabu} (\tsc{Celebic}) 			\\
\ili{Soboyo}	&	\hp{*}fiak	&	\tsc{Taliabu} (\tsc{Celebic}) 			\\
		\lspbottomrule
	\end{tabular}
			\begin{tablenotes}\tnsize
				\item[†]	{Replacing \citet{DonohueDenham2009} listing \ili{Buru} \it{fua-t}.}
				\item[‡]	{Corrected from Tai-\ili{Kadai}.}
			\end{tablenotes}
	\end{threeparttable}
\end{table}
		\ssscp{Lexical taboos require synonyms}{02-03-02-02}
\citet{Simons1982} has written about various kinds of lexical
taboos in Austronesian languages
and pointed out that tabooing assumes
and requires the notion of synonyms -- multiple words
and strategies to talk about the thing you aren’t permitted to
say the normal way.
\citet{GrimesMaryott1994} have written about lexical taboos
and named avoidance speech registers in Austronesian languages,
many of which are associated with hunting and fishing taboos.
\citet{Blust2012b} notes that cuscus are an important animal in the region.
Grimes can verify that from personal experience with multiple
ways of hunting cuscus on \ili{Buru},
trading the dried meat,
and folktales (see \frf{02-21}).

\begin{figure}
	\centering
	\frame{\includegraphics[width=0.375\columnwidth]{02-21.jpg}}
	\fcp{Cuscus meat and broth}{02-21}
\end{figure}

\citet{Grimes1991c,Grimes2000a} gives multiple examples of taboo-generated
synonyms on \ili{Buru} to avoid saying the name of the item being hunted,
or an item considered potentially dangerous.
For example:

 \begin{itemize}
	\item \ili{Buru} \it{uran} ‘shrimp,
				lobster, freshwater crayfish’ (< PMP *qudaŋ ‘shrimp,
				crayfish, lobster’) is more often referred to in the mountainous
				interior as \it{sehe}, which is otherwise a verb ‘reverse,
				back up, withdraw’ associated with the movement of such creatures.
	\item Deer, not native to \ili{Buru} but often hunted,
				are sometimes called \it{menʤaŋan} ‘deer’ (Javanese loan,
				reflecting contact issues), but more often referred to as \it{wadun}
				‘back of neck’, the part of a deer prominently in view when it is feeding head down.
				Loans associated with PWMP *uRsa ‘deer’ (\ili{Malay} \it{rusa}) found
				in many other languages of eastern Indonesia are not known on \ili{Buru}.
	\item The \ili{Rana}
				and Wae Sama dialects of \ili{Buru} refer to mosquitos as \it{seŋet}
				(< PMP *səŋət ‘to sting or bite,
				of an insect’), whereas the \ili{Masarete} dialect has lexicalised
				a descriptive avoidance term \it{inhadat} /ii en-hada-t/ ‘the biting thing’.
	\item Snakes in general on \ili{Buru} are referred to as \it{inewet}
				/ii en-newe-t/ ‘the living thing’,
				and reticulating pythons (\it{Malayopython reticulatus}) are
				called \it{isaleu} /ii saa leu/ ‘the thing that goes before’.
				There are other descriptive terms for specific kinds of snakes,
				with no apparent traces of PMP *hulaR or *nipay.
				Grimes has seen pythons in the \ili{Buru} jungle up to 55cm in circumference,
				and also seen normally fearsome hunter-warriors fleeing them in panic.
				This reflects the perceived danger associated with many lexical taboos.
	\item The elusive civet cat (\it{Viverra tangalunga}
				and \it{Paradoxurus hermaphroditus}, both said to be introduced
				into Maluku from the west,
				see \frf{02-22}), is also feared as a territorial spirit near
				the abandoned mountain village of Ena Biloro.
				Elsewhere it is hunted and eaten.
				The civet cat is called \it{kamyoho} in \ili{Masarete},
				and \it{yoho} in the \ili{Rana} dialect,
				but \it{ŋewaet} in the related \ili{Lisela} language,
				which parses to /ŋee wae-t/ ‘wet nose’ -- an avoidance term.
				Reflexes of PMP *musaŋ ‘civet cat
				and similar’ are not known on \ili{Buru}.\footnote{
						On \ili{Semau} Island
						near Timor, \ili{Helong} uses an avoidance term \it{mat-mitaŋ} ‘black eyes’ to
						refer to the civet cat,
						from the raccoon-like dark fur around the eyes,
						its nocturnal behaviour, and ability to see at night.}
 \end{itemize}

\begin{figure}
	\centering
	\includegraphics[width=0.5\columnwidth]{02-22.jpg}
	\fcp{Asian palm civet}{02-22}
\end{figure}

%\begin{figure}
  %\subfigure[Cuscus meat and broth]{\label{fig:02-21}
    %\includegraphics[width=.45\textwidth]{02-21.jpg}}\hspace{5mm}
  %\subfigure[Asian palm civet]{\label{fig:02-22}
    %\includegraphics[width=.45\textwidth]{02-22.jpg}}
%\caption{Native mammals of Wallacea}
%\end{figure}

These are all examples from \ili{Buru}
and the related \ili{Lisela} language spoken on the same island -- and there are more.
The basic idea of synonyms generated through taboos
and avoidance is repeated in languages throughout the region.
No cognates with *kandoRa or *mans(aə)r have been encountered
in \ili{Buru} or \ili{Lisela} for ‘cuscus’,
although \it{mate} ‘marsupial, cuscus’ has been reported for distantly related \ili{Ambelau}.
And \it{maren} ‘cuscus’ has been reported for recently extinct
\ili{Kayeli}, an \tsc{\ili{Ambon}-Seram} language formerly spoken in eastern \ili{Buru}.
(Note the \ili{Ambelau} and \ili{Kayeli} forms for ‘cuscus’,
and many like them in other languages,
is cognate with *mansər,
which \citet{Blust1982,Blust2012b} glosses ‘bandicoot’, not ‘cuscus’.
Bandicoots are not found as far west as \ili{Ambelau}.) Other competing
terms are also reconstructed for ‘cuscus’,
such as Proto-Celebic *kuse ‘bear cuscus’
and *taŋali ‘dwarf cuscus’,
Proto-\ili{Gorontalo}-\ili{Mongondow} *kusoy ‘cuscus’,
Proto-North Halmahera *kusor ‘cuscus’,
Proto-Nuclear-\ili{Meto} *ʔ-mauka-ʔ ‘cuscus’ \citet{Edwards2021},
and several other reconstructions across Sulawesi,
Halmahera and west Papua.
We return to this issue in \crf{ch:Mar}.

		\ssscp{Parallelisms and doublets}{02-03-02-03}
Throughout the region in all of our Wallacean subgroups,
both parallelisms and doublets are found to varying degrees,
not only in poetic ritual language,
but also in everyday speech
(\citealp{Fox1974,Fox1988,Fox2014,Fox2016}, \citealp[15--31]{GrimesTherikGrimesJacob1997}, \citealp[162--164]{TaberTaber2015}).
Parallelisms tend to be clausal,
while doublets (or word pairs, as illustrated in the next paragraph)
can be hyphenated or joined in a
unitary phrase by a conjunction.
And the two can work in concert,
as illustrated by the translation of part of a \ili{Molo} \it{onen} ‘prayer’ in
\citet{GrimesBIP}: \it{“Lord,
when we encounter difficulty
and suffering, abuse and grumbling, you make us cherish love and inner peace.
Lord, when we encounter abuse
and grumbling, you make us cherish forgiveness.”}
\citet[379--385]{Edwards2020}
discusses the sophisticated interplay of metathesis involved with \ili{Amarasi} parallelisms.
The use of synonyms in parallelisms
and doublets throughout the region are part of what makes language
beautiful, and speakers who use them eloquent.

Whether looking at parallelisms or doublets,
most often similar categories of words or near synonyms are treated
as \it{synonymous} (same meaning) when used together.
Antonyms or \it{counterparts} often reflect whole categories:
mother-father = parents; night-day = around the clock; ascend-descend
= go all over; plate-fork = eating utensils; pigs-chickens = domesticated animals.
And sometimes the pair is \it{elaborative},
fleshing out the picture slightly,
but that is the rarest type.
It is quite common for verbs that have slightly different meaning
in isolation to be used synonymously,
such as \ili{Kotos} \ili{Amarasi} \it{t-kumain ma t-humaʔmoe} ‘we laugh
and we smile’, \it{n-poi ma n-boor} ‘it appears and bursts forth’.
Sometimes these synonymous pairs are easiest to gloss simply
as ‘cut1-cut2’, or ‘carry1-carry2’.
Fox (various) has written extensively about additional figurative
or symbolic meanings when used in ritual language registers (\it{bini})
in the languages of Rote.
An example from nearby Timor,
in the \ili{Kotos} \ili{Amarasi} \it{aʔa sramat} genre,
the phrase \it{mutiʔ ma mnatuʔ} (lit: ‘silver
and gold’) is conventionally used to refer to ‘kings
and nobles’, whereas \it{koro-manu} (lit: bird-chicken) is conventionally
used to refer to ‘populace, commoners’ \citep{GrimesBIP}.
\citet[28]{GrimesTherikGrimesJacob1997} observe “In some pairs,
the meaning of one unit may not be recoverable.
There are several instances where the only occurrence of a word
in the entire available data corpus is in doublet form,
and native speakers \it{cannot} give a meaning for the word in
isolation.” These are more commonly the second member of a doublet
and are often identifiable as being from neighbouring languages or dialects.
\citet{Valeri1990} calls this “the retreat of the semantics”
in central Maluku, in favour of celebrating links
and alliances with other languages
and ethnic groups. \citet[30]{GrimesTherikGrimesJacob1997}
summarise, “The skilful use of parallelisms is at once
a link to the ancestors,
a link expressing solidarity with those of other dialects or
languages with whom one is allied,
and a social link with those who hear
and appreciate the skilful art of the poet.” These patterns of
using doublets and parallelisms to celebrate synonyms,
and taking synonyms from other dialects or languages are well
documented for many languages
and subgroups in the region.
Since the use of parallelisms is also documented in the Philippines
and Borneo, there is no reason to assume it is a recent development
restricted to this region.
\citet{GrimesBIP} also describes similar patterns found in Papuan
languages, such as the TAP language \ili{Klon} on \ili{Alor}.


		\ssscp{Different patterns of contact with people {\&} environment}{02-03-02-04}
\ili{Sumbawa}, Flores, Sumba, \ili{Alor},
Timor, Tanimbar, Seram and \ili{Buru} are all islands ranging in size
from 2,000 square kilometres to over 30,000 square kilometres.
Many of them have significant
and very rugged mountains, and even micro-climates.
On \ili{Buru} the word for ‘mountain’ at first appears to reflect dialect
or language differences with \ili{Rana}
and \ili{Lisela} \it{fude} ‘mountain’ in the north,
and \ili{Masarete} \it{fuka} ‘island,
mountain’ in the south.
In fact, it is a reflection of different geology on the north and south coasts.
The island has a sharp geological southern uplift with cliffs
and tall mountains rising strait from the sea,
and \it{fuka} refers to mountains made of uplifted coral (\frf{02-23} and \frf{02-24}).
The north coast has a narrow coastal plain
and gentler hills composed of soil (see \frf{02-25});
these are \it{fude}.

\begin{figure}
  \subfigure[\it{fuka}: uplifted  coral in south \ili{Buru}]{\label{fig:02-23}
    \frame{
\includegraphics[width=0.31\textwidth]{02-23.jpg}}}\hspace{2mm}
  \subfigure[\it{fuka}: uplifted  coral in south \ili{Buru}]{  \label{fig:02-24}
    \frame{\includegraphics[width=0.31\textwidth]{02-24.jpg}}}\hspace{2mm}
  \subfigure[\it{fude}: mountains of soil in north \ili{Buru}]{\label{fig:02-25}
    \frame{\includegraphics[width=0.31\textwidth]{02-25.jpg}}}
\fcp{Hills and mountains in Buru}{02-23--02-23}
\end{figure}

In addition to synonyms created by taboos described above,
and different environmental features,
there are also dialectal “synonyms” created by different patterns of contact.
For example, \ili{Lisela} in the north of \ili{Buru} (a \tsc{Sula-Buru} language)
has had intense contact over centuries with the recently extinct
\ili{Kayeli} language spoken in the eastern part of \ili{Buru} (an \tsc{\ili{Ambon}-Seram} language).
And for several centuries the king of \ili{Lisela} was subordinate
to and paid tribute to the king of \ili{Kayeli}.
Additionally, there are around 12,000 immigrants from Sula island
to the north, settled along the north coast of \ili{Buru} for 12 or more generations
\citep{Grimes1991c,Grimes2000a,Grimes2010b}.
They have no impact on the south.

PMP *p = \it{p} in \ili{Buru}
and \ili{Lisela} native vocabulary,
as in PMP *pulut ‘birdlime,
sticky sap of breadfruit’ becoming \ili{Buru} \it{polo-n}
and \ili{Lisela} \it{pulu-n} ‘glue,
sticky substance, starchy paste’;
and PMP *sapa ‘what?’ becoming \ili{Buru}/\ili{Lisela} \it{sapa-n} ‘what?’.
But in \ili{Kayeli} PMP *p > \it{h}.
So PMP *paRih ‘skate,
stingray’ becomes \ili{Buru} \it{pahi} ‘stingray’
and \ili{Kayeli} \it{hali} ‘stingray’.
PMP *pajay ‘rice’ becomes \ili{Buru} \it{pala} ‘rice’
and \ili{Kayeli} \it{hala} ‘rice’.
PMP *qapuR ‘lime, calcium’ becomes \ili{Buru} \it{apu} ‘lime,
chalk, calcium’ and \ili{Kayeli} \it{ahul} ‘lime, chalk, calcium’.
So \ili{Lisela} \it{hali} ‘stingray’,
\it{hala} ‘rice’, and \it{ahul} ‘lime, calcium’ identify themselves
as \ili{Kayeli} loans, rather than inherited reflexes of the PMP forms.
On a mountainous island of nearly 10,000 square kilometres,
people in the south of \ili{Buru} know \it{apu}, but not \it{ahul}.
People in the north of \ili{Buru} know \it{ahul}, but not \it{apu}.
The \ili{Rana} dialect in the middle celebrates its links through marriage
both to the north and south.
So it is not surprising that the sung ballad of Raja Bau with
195 stanzas in poetic parallelisms uses both \it{apu}
and \it{ahul} as synonyms.

Like many Austronesian languages,
\ili{Buru} has phrase level genitive enclitics
(at the end of the word
to which it attaches).
So \it{(geba) faha=n} ‘(person) hand/arm-\tsc{3sg.gen}’ (< PMP
*qabaRa ‘shoulder’), \it{faha=ŋ} ‘my arm’, \it{faha=m} ‘your arm’,
\it{pani=n} ‘its wing’ (< PMP *panid ‘wing’),
\it{aa=n} ‘his jaw’ (< PMP *qazay ‘chin’), etc.
Sula, an island to the north with thousands of immigrants on
the north coast of \ili{Buru},
many of who say their line has been there for 12+ generations,
has preposed genitives, as in \it{n-aya} ‘his jaw’.
Both \ili{Buru} and Sula (\ili{Sanana}) use their genitives
in vocative constructions (direct address).
So a \ili{Buru} person (\ili{Rana} dialect) addresses his parents as \it{ama=ŋ}
‘my father’ and \it{ina=ŋ} ‘my mother’ (< PMP *ama ‘father’; *ina ‘mother’),
whereas a Sula person addresses his parents as \it{ŋ-ama} ‘my
father’ and \it{ŋ-ina} ‘my mother’.
The \ili{Lisela} vocatives are the same as the Sula forms,
indicating they are Sula loans.
The languages in the Sula Islands are the only ones with preposed
genitives with which \ili{Lisela} has had significant contact.

There are also many non-cognate forms that differ in different
dialects, some triggered by different patterns of contact,
some by tabooing, and the motivation for others remains unclear.
So ‘cuscus’ is \it{tonal} in the \ili{Masarete} dialect of \ili{Buru} (the
final \it{l{\#}} marks this word as unusual in its structure,
suggesting it is a loan from a south\ili{east} Sulawesi source -- compare
\ili{Kambowa} \it{tonali} ‘dwarf cuscus’ [\it{Strigocuscus celebensis}]),
\it{lemet} in the Wae Sama dialect,
and \it{blafen} ‘cuscus (generic)’
in the \ili{Lisela} language.
‘Mosquito’ is \it{inhadat} in the \ili{Masarete} dialect (a taboo synonym),
\it{seŋet} in the \ili{Rana}
and Wae Sama dialects (PMP retention with a shift in meaning),
and \it{mimaun} in the \ili{Lisela} language.
We find similar variation for skin,
lips, wing, megapode, thatch,
fish trap, wave, sleep, wake up, hungry,
cold/flu, defecate, good, many,
husband, wife, cloth, machete,
medicine, quick, some verbs of speaking,
some prepositions, and many other things.
Why would they need competing terms to talk about the same thing
on the same island? The simple fact is,
they do have competing terms for many things
and for a variety of reasons.

Some languages may appear to have ``competing words'' that actually
refer to different species or varieties.
For example, \ili{Lisela} (\tsc{Sula-Buru}) has \it{blafen} ‘cuscus
(generic)’ and \it{ton gawat}
‘referring to the spotted cuscus (\it{Spilocuscus maculatus})’.

Similar variation is found in other larger languages scattered
across mountains and coasts, and having different areas with different contact
patterns, such as \ili{Alune}, \ili{Wemale},
\ili{Manggarai}, \ili{Sika}, \ili{Lamaholot}, \ili{Kambera},
the \ili{Meto} languages, the Rote languages,
\ili{Tetun}, \ili{Mambae}, and \ili{Yamdena}.

So, when \citet[268]{Blust2012b} says, “we must ask why a language
would have synonyms for ‘cuscus’”,
the response is fairly simple,
“There are many reasons,
and they are all very well documented in the region.” Appealing
to the situation in Hawai‘i where the Austronesians encountered
new flora and fauna in a context of no prior human habitation,
and where seeking loan sources is futile,
\citet[273]{Blust2012b} boldly declares about the dispute over
marsupial terms in eastern Indonesia that,
“there was also no reason for \tsc{An} immigrants to borrow the
name of the cuscus rather than invent one from scratch.”

Yet the Hawai‘i situation is not at all parallel to the situation
in our target region.
\ili{As} we have shown (\srf{02-01}),
there is evidence of over 40,000 years of human habitation.
The Austronesian-speaking peoples did not enter a previously uninhabited region.
In \trf{02-Features} we gave several examples in which Austronesian
languages borrowed Papuan ways of saying things when they already
had perfectly good ways of saying them in their existing communication
framework, such as ‘hair {\tl} leaf of head’,
‘liver {\tl} insides’, ‘banana’,
the preposed possessor, alienable-inalienable distinctions,
clause final negation, and so forth.
Why invent a word,
when other people that you interact with,
who are not enemies,
already have names for it? Coining a new word “from scratch”
is usually quite rare.


		\ssscp{Receptive to outsiders and outside ideas}{02-03-02-05}
In addition to abundant evidence of Austronesians borrowing
words, metaphors, and even grammatical structures,
but they were also receptive to ``outsiders''
and not ``closed'' in other aspects of life.\footnote{
		The statements here about Austronesian-speaking peoples
		being receptive to outside influence also apply to those
		speaking Papuan languages (e.g. \citealp{Schapper2011-BunakHistory}).}
For example,
\citet{Sahlins1985,Sahlins2008,Fox1995,Grimes2006,Haegerdal2008}
are some of those who have written about the
widespread pattern of accepting ‘outsider kings’
or ‘stranger-kings’ in Austronesian societies.
Given the abundant evidence of people speaking Austronesian languages
borrowing words, metaphors, structures, and being receptive in other ways to outsiders
and their ideas, it is extremely likely that the names for new
flora and fauna encountered as they spread through the target region
were borrowed from other existing terms readily available in
the region, rather than invented from scratch.
The severe under-documentation of languages in the region,
combined with multiple sources of evidence that Papuan languages were once
spoken much more widely in the past (some wiped out by volcanos),
means we may never be able to identify a specific source language.
But that does not make borrowing any less likely.
Arguments in the absence of data do not move us forward.
So we are extrapolating from patterns for which we do have data.
But as shown in \crf{ch:Mar},
there are also Austronesian languages west of the Blust line
which have words of (incrementally) similar form which mean ‘cuscus’,
so a Papuan source is not the only possible explanation.

In response to \citet{Schapper2011}
and \citeauthor{DonohueGrimes2008}' (\citeyear{DonohueGrimes2008}) suggestions that Papuan trade
networks provide a mechanism for the diffusion of words as well
as goods,\footnote{
		Such as those responsible for the human introduction
		of the cuscus from the New Guinea region into western Maluku
		and the Lesser Sundas,
		and bananas, sago, taro and sugarcane as far as western Flores.}
 \citet[274]{Blust2012b} says,
“Since trade is a universal human phenomenon,
it clearly existed in the Papuan world on a local scale before
AN contact, but there is no evidence that it was pan-archipelagic,
or that any Papuan-speaking group had seafaring capabilities
at all comparable to those of AN speakers when they arrived in
the area with the outrigger canoe complex.”

On the “seafaring capabilities”,
we have already shown that there is growing archaeological evidence
that allowed deep water crossings
and open ocean fishing in the region long pre-dating the arrival
of the Austronesian-speaking peoples
\citep{Spriggs2011,OConnor2015,Bellwood2019ed}.
Numerous scholars have noted that the westward spread from the
New Guinea region of bananas,
sugarcane, taro, sago and breadfruit strongly implies a significant
seafaring capability, again predating the Austronesians
\citep{DenhamDonohue2009,DonohueDenham2009,DonohueDenham2010,Schapper2015,
Klamer2019,GrimesCIP}. There is no need to argue whether or not
it was “comparable” to Austronesian capabilities,
since whatever the technology was,
was clearly adequate for the dispersals we are discussing.
Add to that, the \ili{east}
and west prevailing winds shift back
and forth on a yearly seasonal basis,
dominate the weather, agricultural cycles,
and hunting cycles throughout the region,
and are a focus of all boatmen across the islands,
regardless of size or type of boat,
including dugout canoes without outriggers.

But the kind of dispersal of food crops
and the wide distribution of fairly stable terms like {\#}muku
‘banana’ can -- and apparently did -- happen through simple local
networks of contact, without having to appeal to a “pan-archipelagic” trade network,
or to a single Papuan trade language or lingua franca.
Things for which there is,
as yet, little evidence.
(In fact there is evidence of other long-standing trade networks
that contribute to the complexity in the region,
described in \srf{02-02-01}, but not all of them predate the
arrival of Austronesian-speaking peoples.)


\largerpage
\scp{Complexity in the language documentation}{02-04}
A number of issues in documenting,
transcribing and analysing the languages have also contributed
to muddying the waters,
making it harder to gain clarity in unravelling the prehistory
of the languages in the region.

	\sscp{Size of sample for establishing patterns}{02-04-01}
We have already noted that the Austronesian languages in our
target region are severely under documented,
with very few published grammars or dictionaries,
a few short grammar
and phonology sketches, and a few additional papers on other
languages about limited topics (e.g. negation, subjecthood, possession).
While for several languages we have access to more recent,
more reliable, and more complete data,
much of the earlier comparative discussion is based on word lists,
many of which are very limited (200 items or less),
and of varying quality and reliability.
The sample is often simply too small to make definitive statements
about regular sound correspondences or to have confidence in
some of the glosses.

\citet[8]{Klamer2019} quantifies it this way:

\begin{quote}
	“The data sets on which earlier classifications of languages
	in this region have been based are very limited and narrow.
	For example, of the 208 languages that Glottolog
	(Hammarström, Forkel, {\&} Haspelmath, n.d.) [Glottolog 3.0] lists as spoken
	in eastern Indonesia, less than 50{\%} (102) are represented in the major reference
	work on Austronesian reconstructions,
	the Austronesian Comparative Dictionary
	(ACD; Blust {\&} Trussel, n.d. [\citet{BlustTrussel2020}]).
	Of those, about 50{\%} (52) features in the database with less
	than 10 words, and 25 of them with only one (!) word.
	[{\ldots}] Clearly, affiliations based on so little data collected
	from such a vast insular region comprising thousands of islands
	leave a lot of room for uncertainties.” \citep[8]{Klamer2019}
\end{quote}

Ten words in a language,
or even a hundred,
is not necessarily going to reveal,
for instance, that PMP *R > \it{l} word initially,
but *R > \it{r} between vowels,
or that PMP *j > \it{l} in some words or environments,
but *j > {\0} in others in the same language.
With such a small data sample we have little basis by which
to gauge the accuracy of the transcription (e.g.
are glottal stops indicated? double vowels? schwa? is orthographic
‹b› fricative [β], implosive [ɓ]?
Is ‹f› bilabial [ɸ]?) or the reliability
and completeness of glosses (e.g.
Is the word for ‘sun’ also the word for ‘day’? Is the word for
‘boat’ generic, or refers only to ‘outrigger canoe’? Does the word given for
‘bird’ also include other winged creatures such as bats?).
Is the word for ‘grandparent’ a reciprocal term that also includes
‘grandchild’? Is the word glossed ‘brother’,
from the perspective of a same sex sibling,
opposite sex sibling, or either?

\ili{As} seen in the following chapters,
some comparative linguists are willing to draw conclusions
and make subgrouping arguments on the basis of historical medial
*NC consonant clusters, such as *mp, *mb,
*nt, *nd, *ns, *nz,
*ŋd, *md, *ŋk, *ŋg.
But the reality in the data available to us in our target region
shows a variety of problems that should encourage caution.

\begin{enumerate}[label=(\alph*)]
	\item	The number of proto forms for which we have data in which
				such sequences occur is quite low,
				often an occurrence of only one or none.
				This is true even for fairly well documented languages.
	\item	The number of words which reflect such medial clusters
				found in enough daughter languages to see a comparative pattern
				is also quite low to non-existent.
	\item	Some languages reflect the presence of historical clusters
				and others do not (e.g. *upu vs.
				*umpu ‘grandparent/child’, *ŋisi vs. *ŋinsi ‘tooth’).
				To illustrate a number of issues under discussion:
		\begin{itemize}
			\item In \citet{BlustTrussel2020}, \ili{Buru} \it{upu}
						is assigned to PMP *umpu ‘grandparent/ grandchild (recipr.); ancestor’.
			\item But \it{upu} is the form for the \ili{Lisela} language.
						The \ili{Buru} form is actually \it{opo}, and \it{opo=n}.
			\item \ili{Buru}
						and \ili{Lisela} sound correspondences are both *p = \it{p}, *mp > \it{b}.
						These are very consistent (except in a few loan words \ili{Lisela}
						has borrowed from \ili{Kayeli} which reflect *p > \it{h}).
			\item The expected form from *umpu for \ili{Buru} would be an unattested *[obo].
						So we can say with a high degree of confidence that \ili{Buru} \it{opo}
						is from a parent form of *upu, not *umpu.
						(Some nearby languages do reflect *umpu,
						but \ili{Buru} and many other languages in the same region reflect
						*upu; still others are not clear.)
		\end{itemize}
	\item	A single historical cluster may reflect three or four different
				patterns in the languages (see \citealp{Grimes1991a}).
				For example, *ma-diŋdiŋ ‘cold’ may be retained in different languages
				as \it{didi, diŋin, rini, marindin, marihi,
				ridi, riki}, or \it{lirin,} reflecting several different responses
				to the medial cluster.
	\item	A single language may reflect more than one response.
				For example, \ili{Kei} Kecil has both \it{ŋa-ridin}
				‘cold’ and \it{na-b-rinin} ‘chills’.
	\item	A single language may reflect an *NC cluster in one word,
				and a medial *C without the nasal in another word.
				For example PMP *upu yields the expected \ili{Buru} \it{opo,
				opo=n} ‘1) grandparent-grandchild (reciprocal term); 2) master, lord’.
				Whereas the related PMP *ta umpu ‘grandparent/grandchild (referential
				term used reciprocally)’ yields the expected \ili{Buru} form \it{tobo=n}
				‘person responsible, master, lord, one with custodial rights, owner, sir, host’.
				So care needs to be taken not to assume all reconstructed medial
				*NC clusters behave the same way in the same language for all words.
				This issue is troublesome for the whole idea of regular sound
				correspondences, unless one treats the medial *NC cluster as optional for many
				words in many languages,
				and checks each item separately.
	\item	\ili{As} above,
				two words in a language may reflect two different patterns for
				the same historical cluster.
				For example, the reflexes for *ns in *ŋinsi (local gloss: ‘tooth’)
				show a different pattern than for *mans(aə)r (local gloss: cuscus).
				There may be no other words found in that language reflecting
				*ns in a proto form.
				An occurrence of one is not a pattern
				and does not qualify as a regular sound correspondence.	
\end{enumerate}

We tend towards caution on making subgrouping claims based on
rare or elusive historical *CC clusters -- unless there is actually
sufficient data to establish a pattern
and substantiate a claim.
We prefer to see \it{a minimum} of two to three examples showing
the same pattern before allowing that it might reflect a regular sound correspondence.
We remain sceptical about claims of sound correspondences for
rare or elusive data reflecting *r,
*nj, *nz, *ns, *ŋk,
*ŋg, for example, since we ourselves can find very few unambiguous reflexes of these.
Even more rarely do we find enough reliable examples in related
languages to base claims about subgrouping on them.

What a small sample of ten words or so sometimes may allow,
is a glimpse of how it is similar to,
or differs from other languages in the same local region or node of classification.
The small sample may (or may not),
be enough to reflect a diagnostic feature for that node.

To take one example,
\citet[129]{Collins1986} classified \ili{Seti},
spoken in central-eastern Seram,
with \ili{Masiwang} and \ili{Bobot} partly on the basis of shared *ŋ > \it{k},
though he noted: “In \ili{Seti} the reconstruction is a problem because
so few cognates containing *ŋ occur.
Two display \it{n} while one shows \it{k}.” However,
examination of more data (to which Collins did not have access)
shows that the one word with *ŋ > \it{k} (*naŋuy > \it{nak} ‘swim’)
truly is an exception with ten other examples of *ŋ > \it{n}.
(See \srf{10-03-04}, particularly example
\ref{ex:10-59} for more discussion of \ili{Seti}).
\citet{Collins1986} classification was reasonable based on the
limited data available to him,
but it is illustrative of the problems
faced with many languages across the region.



	\sscp{Labelling of language names}{02-04-02}
Language names may be mislabelled,
misleading, or are actually the names of villages.
For example, many readers are unaware that \citet{Collins1983}
tends to label his data according to the village name in which
a word list was taken or from which the consultant came,
and not necessarily the language name.
So some readers are unaware that several data sets are from different
varieties of the same language such as \ili{Alune} or \ili{Wemale}.
Only a careful reading may help sort that out.
Even those reasonably familiar with the wider region may be easily confused.

Similarly \citet{BlustTrussel2020} list data from “Tifu”.
There is no Tifu language,
and it should not be treated as a separate witness.
The data are clearly from the \ili{Buru} language.
Tifu is a small village on the south coast of \ili{Buru} where Dutch
UZV (\it{Utrechtse Zendings Vereniging}) missionaries Hendricks and Schut resided.
They both published on the \ili{Buru} language
and culture, and were the main sources for
\citeauthor{Stresemann1927}'s (\citeyear{Stresemann1927}) \ili{Buru} data.

In several of our subgroups,
varieties that share the same language name in the literature
may have patterned differences in their sound correspondences
that are important for historical-comparative work.
This is true for \ili{Alune},
\ili{Wemale}, \ili{Kei}, “Rote”, “\ili{Meto}”,
“\ili{Manggarai}”, and other large or geographically widespread languages.

In some cases, we have been unable to confidently identify which
language a particular set of data represents or its geographic
origin due to problems in labelling.
Thus, \citet{Tauern1931} has short grammar sketches of a number
of languages of Seram.
One of these is labelled “Uhei-Kaẖlakim”
and another is labelled “\ili{Liambata}”.
Based on comparison with data in other sources,
these are probably varieties of \ili{Seti} (\srf{10-03-04}),
but we are not completely confident.
\citet{Stresemann1927} also has “\ili{Liambata}” data which is different
to that in \citet{Tauern1931},
and, to further complicate matters,
\citet{Tauern1931} also refers to his “\ili{Liambata}” as “\ili{Bonfia}”,
which is the Dutch-era label for \ili{Masiwang} [bnf].
Comparison with sources which unambiguously present \ili{Masiwang} data
(e.g. \citealp{LeCocqdArmandville1901} and
\citealp{Collins1986})
shows that the “\ili{Liambata}” data in \citet{Tauern1931}
is definitely \it{not} \ili{Masiwang},
but it is easy to see how confusion
and mis-identification can arise.


	\sscp{Reliability of the sources}{02-04-03}
Comparative linguists’ handling of the data can be only as good
as the sources of their data,
some of which are already secondary sources,
and some of the primary sources may be mislabelled,
under-differentiated, actually represent another language or
variety, be poorly transcribed, or misglossed.

For example, \citet{BlustTrussel2020} list a large quantity of
data under the label of “Buruese”.\footnote{
		The label “Buruese”
		is not from the \ili{Buru} langauge, it is not Indonesian, and not even good \ili{English}.
		Old Dutch sources (24 publications by Hendriks,
		Jansen, Jellesma, Schut, van der Miesen, Wilken,
		Willer) refer to the island,
		people and language as “Boeroe”
		and “Boeroetaal”, with one instance of “Burusch”.
		Devin also used “\ili{Buru}” for his unpublished dictionary.
		“\ili{Buru}” suffices to identify the language,
		and is much preferred.}
 Much of that data is from an unpublished dictionary by Ch.
Devin who did not distinguish different varieties of \ili{Buru} or
the \ili{Lisela} language, but tended to lump everything together.
Devin’s main informant, Bapak Temi,
was from Banu Lalet,
near the headwaters where the \ili{Lisela},
\ili{Rana}, Wae Sama, and \ili{Masarete} varieties converge.
However, several forms cited as “Buruese” are identifiable by
form or meaning as actually being from the related \ili{Lisela} [ISO:
lcl] language -- which has some patterned differences with the
\ili{Buru} language [ISO: mhs] in both the phonology (e.g.
different handling of PMP *u,
different handling of PMP *p,
sometimes a different vowel in antepenultimate syllable,
sometimes different CC onsets; \citet[37--38]{Grimes1991c}
and the lexicon (lexical similarity of \ili{Lisela} with several \ili{Buru}
dialects averages only 60--63{\%}).
Again, this reflects on the source,
but ultimately muddies the historical-comparative work.

From 1895--1939 several versions of an approximately 1,000-item
word list known as the “Holle lists” were prepared
and sent throughout what is now the Indonesian archipelago.
\citet[17]{Stokhof1980a} notes, “It was planned that the lists
would be completed by Dutch civil servants,
officers, missionaries and ‘intelligente Inlanders’ such as village
heads, merchants and teachers, in practice all people with a reasonable
knowledge of Dutch.” A number of scholars have noted the unreliable
labelling of language varieties or transcriptions or glosses
in some Holle lists.
A common problem is that those who filled out the lists were
often from entirely different regions,
different sound systems, and different semantic ranges,
untrained in linguistics, showing goodwill in trying to help
collect the desired data.
We are grateful for that,
but some word lists and whole sections of the data can be notoriously unreliable.
\citet[10--11]{Grimes1991c} critiques several Holle lists associated
with \ili{Buru} and \ili{Kayeli}.
For example, the Holle list labelled “\ili{Hukumina}” \citep[63ff]{Stokhof1982b}
is not just the recently extinct \ili{Hukumina} language from western
\ili{Buru}, but also contains a mixture of the \ili{Lisela}
and \ili{Kayeli} languages from a village named \ili{Hukumina} near \ili{Kayeli}
village in eastern \ili{Buru} \citep{Grimes2010b}.


	\sscp{Issues with glosses}{02-04-04}
Some elusive terms or glosses are central to debates about Austronesian prehistory.
For example, \citet{Fox2011} points out that an 1889 Dutch source
for certain \ili{Ngaju} Dayak sibling terms has not been able to be
verified by other anthropologists (James Fox pers. comm. October 2020),
nor is it known by native speakers.
Verifying the term (or not) is important since it appears to
be an anomaly in the distribution of cognate forms,
most of which are found far to the \ili{east} in Linguistic Wallacea.
A report from nearby Uut Danum shows possible cognates restricted
to formal speech and oral epics.
The terms in Borneo are elusive at best,
and may have been present in certain language(s) in the past,
but are no longer current.
Yet the existence and accurate glossing of such sibling terms
denoting relative age and opposite sex are
critical to understanding early Austronesian
social structure and transformations
in different regions (see \citealp{Fox2015}).

Difficulties with glosses impinge on precisely identifying \it{flora
and fauna} for comparative purposes.
For instance, word lists with ‘grass’ often list the first answer
given, but the investigator may or may not explore further to see if
it is a generic term,
or if there are different terms for saw-toothed vs.
smooth edged varieties -- which often have different historical etyma.
Is the word given for ‘bamboo’ a generic term,
or does it refer to a certain species of bamboo? Word lists with
‘uncle’ often do not distinguish between ‘mother’s brother’
and ‘father’s brother’, which are not distinguished in the languages
of elicitation, but are consistently distinguished in the social structures
and lexicons of our target region.
Investigators taking word lists may assume sibling terms have
only ‘elder sibling’ and ‘younger sibling’ distinctions
and fail to check if there are any ‘opposite sex sibling’,
‘older sibling (same sex)’,
or ‘younger sibling (same sex)’ terms which may not exist where
the investigator is from or in the languages of elicitation,
but are common in the target region.

PMP *duRi is reconstructed meaning ‘thorn’,
which is the overwhelming gloss in languages in the northern
and western Malayo-Polynesian regions.
East of the Blust line the dominant gloss shifts to ‘bone’
and continues out into the Pacific (see \frf{15-02}).
Yet some investigators in eastern Indonesia who have Austronesian
awareness about the etymon meaning ‘thorn’,
have glossed reflexes of *duRi only as ‘fish bone’,
whereas investigators without such bias or awareness gloss them
simply as ‘bone’ in the same language.
It is the word given when working through
lists of human body parts or pig anatomy.

\ili{As} mentioned before, our target region sits between two biogeographical
zones, and also has flora
and fauna endemic to the region.
While there are resources on the flora
and fauna of the New Guinea region (Sahulland),
and of South\ili{east} Asia (Sundaland),
there are few focused on Wallacea.

In discussing discrepancies between the distribution of bandicoots
(the rodents) and the distribution of words glossed as ‘bandicoot’,
\citet[262--263]{Schapper2011} observes, “Six languages with claimed
reflexes in this region do not occur in the area where bandicoots are found.
Furthermore, the terms do not denote ‘bandicoot’.
\citet[110]{Blust2002} makes light of the discrepancy between the meaning
of reflexes and that of his reconstruction,
stating that ‘in some cases reflexes of this form are said to
indicate the phalanger, or cuscus, and it is
not entirely clear whether this is due to
semantic change or to confusion on the part of the lexicographer.’
In fact, with only a single exception,
all reflexes of *mans(aə)r ‘bandicoot’ west of New Guinea denote
‘cuscus’ and not ‘bandicoot’, even in the area where bandicoots are found.”

So Blust appeals to “confusion on the part of the lexicographer”,
yet he also appeals to the opposite when a critical piece of
data supporting his hypothesis is challenged.
\citet[263]{Schapper2011} points out that the \ili{Leti}-\ili{Moa} term \it{mada
{\tl} made} has been incorrectly interpreted as ‘bandicoot’ by
Blust and the correct gloss as ‘cuscus’ has been verified by
\citet{vanEngelenhoven2004}
and that the \it{madar} term reported for \ili{Damer} whose gloss is
translated as ‘marsupial sp.,
bandicoot’ is actually outside the area where bandicoots are found.

\citeauthor{Blust2012b}'s (\citeyear[264]{Blust2012b}) response to this is curious.
Ignoring the information that \citeauthor{vanEngelenhoven2004} verified the \ili{Leti}-\ili{Moa}
gloss as ‘cuscus’, he says, “The one clear discrepancy that remains in this table
is the gloss for \ili{Leti}-\ili{Moa} \it{mada/made},
which \citet{Jonker1932} cites as buidelrat (‘pouch rat’ = bandicoot).
{\ldots} It is odd that Jonker,
who was a very experienced lexicographer
and a scholar well acquainted with the land,
peoples, and languages of the Lesser Sundas,
would commit such an error.” We note that Jonker’s dictionaries
(like most first dictionaries of any language) are full of dozens
of vague entries saying “soort X [soort boom]” ‘a kind of X [a kind of tree]’.
That does not diminish his skill as a lexicographer,
it only points out he was working with many limitations in the
specific identification of flora and fauna.

The debate between \citeauthor{Blust2012b}
and \citeauthor{Schapper2011} illustrates another common methodological problem
of filtering obscure glosses through multiple languages,
and the caution and grace needed in sorting things out.
For example, were the individual Holle word lists done from the
late 1800s to early 1900s taken using Dutch as the language of
elicitation, or \ili{Malay}? Or did some use Dutch
and some use \ili{Malay}? How proficient was the investigator taking
the word list (Dutch investigators using \ili{Malay}; local investigators
using Dutch)? What about proficiency in the regional variety
of \ili{Malay}, which can have significant differences in the vocabulary from
Standard \ili{Malay}? (See examples below.) How proficient in the language
of elicitation (Dutch or \ili{Malay}) was the informant supplying the
data for the word list (some of whom were second language speakers
of the target language).
And then how well do \ili{English} speakers using the materials handle
Dutch or \ili{Malay} glosses from a previous era?

For example, Dutch \it{buidelrat} parses literally to ‘pouch
rat’, so assuming it means some type of ‘marsupial’ is reasonable.
And Dutch \it{rat} means ‘rat’,
so assuming ‘pouch rat’ refers to ‘bandicoot’ is not unreasonable.
But several modern Dutch dictionaries give ‘possum’ (or ‘opossum’)
for \it{buidelrat}.\footnote{
		For example, putting \it{buidelrat} into the Dutch version of Wikipedia redirects
		to opossum (\it{Didelphidae}) -- the family of marsupials found
		in the Americas. (‘Opossums’, Wikipedia.
		nl.wikipedia.org/wiki/ Opossums. Accessed 22 June 2023)}
The Holle list master \citep[50]{Stokhof1980a} equates: item {\#}819
Dutch: \it{buideldier} = \ili{English}: \it{marsupial} = Indonesian: \it{kuskus} ‘cuscus’.
Which language was used for elicitation? The Dutch term could
cover any marsupial, whereas the \ili{Malay} is limited to the \it{Phalanger} ‘cuscus’.
\ili{Ambon} \ili{Malay} has \it{kusu} for ‘cuscus’,
perhaps to disambiguate it from \it{kusu-kusu} ‘kunai grass,
sword grass’ (Indonesian: \it{alaŋ-alaŋ}),
which in rapid speech is pronounced \it{kuskusu}
often with an utterance final decrescendo.
\citet[304f]{vanEkris1864} provides independent evidence via \ili{Malay}
that Dutch \it{buidelrat} is directly associated with ‘cuscus’
during the era when many Dutch sources recorded languages for
which we have no other data.
For the \ili{Kamarian} language of southwest Seram (now said to be
extinct), \citeauthor{vanEkris1864} gives \it{maker} ‘buidelrat’,
\it{make mauriat} ‘gevlekte buidelrat [Ml.
kusu potar]’ (Ml. = \ili{Malay}), \it{make moon} ‘groote buidelrat’.
\ili{As} mentioned above \it{kusu} is the \ili{Ambon} \ili{Malay} term for ‘cuscus’.
And Dutch \it{gevlekte} = ‘spotted’,
so \it{make mauriat} ‘gevlekte buidelrat’ being explicitly identified
in \ili{Malay} as a type of \it{kusu} refers to the ‘spotted cuscus
(\it{Spilocuscus maculatus}; see \frf{02-26})’.
This is fairly clear evidence that the term \it{buidelrat} in
19th century Dutch sources can refer to ‘cuscus’.

\begin{figure}
	\centering
	\includegraphics[width=0.45\columnwidth]{02-26.jpg}
	\fcp{Spotted cuscus \it{(Spilocuscus maculatus)}}{02-26}
\end{figure}

So given the evidence from \citet{vanEkris1864}
and modern Dutch dictionaries,
both \citeauthor{Jonker1932} and \citeauthor{vanEngelenhoven2004} agree that \ili{Leti}
and \ili{Moa} \it{mada/made} refer to ‘cuscus’,
and we do not need to discuss whether Jonker made an error or not.
Similar cross-linguistic confusion can be found with terms referring
to civet cats and bandicoots as well as types of birds,
bats, insects and plants.

To expand on the difficulties in identifying flora
and fauna, some might also characterise both Grimes
and Edwards as “experienced lexicographer[s]” (see Bibliography).
Yet we are happy to acknowledge difficulties in identifying terms
in the field for elusive animals such as the bandicoot
and civet cat, or even more common plants or birds or insects.
There are many variables.
For example: coastal people tend to know very little about the
flora and fauna of the interior or the highlands,
and \it{vice versa} with mountain people,
especially on the larger islands.

The lexicographer can be reasonably fluent or even very fluent
in the local language,
but be talking about an unknown thing,
or limited to verbal descriptions.
For example, \ili{Buru} men in the interior have described to Grimes
a \it{faf bile} that is said to be distinct from long-legged
wild pigs (\it{faf aba}),
distinct from domestic pigs with slightly shorter legs (\it{faf
fena}), distinct from domestic pigs that have gone feral (\it{fafu melaban}),
distinct from wild pigs that have been caught young
and domesticated (\it{fafu emlahat}),
and distinct from the babirusa (\it{faf donit}).
The \it{faf bile} is said to live in the jungle
and “has short legs and fine hair”.
Several conversations with additional hunters who also spend
considerable time in the jungle have not added any new information.
More research is clearly needed.
But when is that likely to happen?

More commonly the comparative researcher has limited fluency
or may even be very fluent in a language of wider communication,
or an official language,
both of which add additional filters
and limitations in both directions.
Many people in remote communities often have no reason to talk
about local flora and fauna in the national language
and often don’t actually know the official language terms.
People in \ili{Buru} know a bit of \ili{Ambon} \ili{Malay}.
People in west Timor know a bit of \ili{Kupang} \ili{Malay}.
The \ili{Ambon} terms and the \ili{Kupang} terms for some of the same types
of plants and insects and fish can be quite different from each other due
to different patterns of contact.
Add to that the problem that most of the standard Indonesian
terms for flora and fauna in major or official dictionaries are biased towards
terms used in the Sumatra-Java region,
and are often oblivious to terms prevalent
in the different regions of eastern Indonesia.

For example, the Indonesian Bible term for ‘cow’ used in Genesis
41:1-4 for pharoah’s dream is \it{ləmbu},
which is not correctly identified as any type of ‘cow’
and sometimes not even as an animal by 70{\%} of 270 tertiary
level theology students and ministers tested in eastern Indonesia by Grimes.
A few recognised it as “an animal used in sacrifices”.
These are not uneducated people,
and they all know the Indonesian word \it{sapi} ‘cow’,
but they live in a region where nobody
uses \it{ləmbu} outside the Indonesian Bible.

Similarly in discussing ways to translate the story of the flood,
the standard Indonesian term for the ‘raven’ released by Noah
in Genesis 8:6--7 is \it{burung gagak}.
In several different remote language communities in eastern Indonesia,
when people were asked to describe what that might be,
there were several similar responses separated by time
and distance and language region saying it might be white or
speckled, and it eats fruit.
Again, these people are not ignorant.
They just don’t know that official language term.

When Grimes was assisting the Ministry of Education with multilingual
education materials in the \ili{Kemak} \ili{Marobo} area of Timor-Leste,
several farmers of both sexes took four days to come up with
the \ili{Kemak} word for ‘ladybug’ for a children’s book selected by
decision-makers in the city.
When they finally recalled it,
it was not a word recognised by most speakers.
A ladybug has no local cultural value
and figures in no folktales.
Several knew the Portuguese loan \it{joaninho} from school books,
because the insect has cultural value in Europe for children’s stories.

But even with the aid of language fluency
and drawings and digital photos the issue is not always clear.
People are used to seeing things in the jungle in colour,
in three dimensions, in their habitats, and in motion.
We are often asking them to interpret something static in two
dimensions out of context,
or from verbal description alone.
Line drawings without colour can also have some drawbacks.
When Grimes showed the black
and white drawing to five men who live in the remote
mountains and spend considerable time hunting,
three different names of doves/pigeons were identified,
discussed and argued over, and a fourth
possibility was named and then discarded.
It was not until he showed the colour version of the drawing (\frf{02-27})
that they eventually reached consensus that it was the \it{man rahe}.
But these visual aids to identifying local flora
and fauna are often unavailable,
or from different biogeographical zones,
or one may be limited to working verbally.


\begin{figure}
	\centering
	\includegraphics[width=0.6\columnwidth]{02-27.jpg}
	\fcp{Emerald dove \it{(Chalcophaps indica)}}{02-27}
\end{figure}

Grimes has also gotten different responses to questions like
“Describe the \it{goli} bird?” Some more experienced hunters
described it as a medium sized meat-eating bird with black wings
that hunts live prey and whose call is \it{ekrikik}.
A more home-body type described it as a seed-and-fruit eating
bird whose call is \it{wio}.
The confidence level is higher for the first one
and whose description probably points to the ‘black-winged
kite’ (\it{Elanus caeruleus hypoleucus Gould}).
But discrepancies mean that further checking is warranted,
and it is unknown when that might happen.

%\begin{figure}
    %\subfigure[]{
		%\includegraphics[width=\textwidth]{02-27a.jpg}
		%}
    %\subfigure[]{
		%\includegraphics[width=\textwidth]{02-27b.jpg}
		%}
	%\fcp{Emerald dove \it{(Chalcophaps indica)}}{02-27}
%\end{figure}

Having experienced many uncertainties in identifying flora
and fauna, and having gotten a few things wrong ourselves,
we are willing to countenance the possibility that some published
sources might not have gotten everything quite right,
particularly when it is anomalous to the rest of the data.
Subgrouping arguments should not hinge on disputed data,
and particularly not on disputed lexical data or their glosses.

	\sscp{Issues with transcription}{02-04-05}
There are many issues relating to problems that arise from the
transcription of words that muddy the picture for comparative-historical purposes.
We discuss here only a sample for illustrative purposes.

Are intervocalic semi-vowels phonemic? phonetic? or simply transcription
variants? For example, words like \it{sawan, awan,
dawe, tewa} are fairly straightforward.
But what about \it{duwe {\tl} due,
ua {\tl} uwa}, \it{leya {\tl} leiya {\tl} lea}? Or \it{tewa {\tl}
teiwa {\tl} taiwa}? This issue is relevant to many languages
in the region, and to understanding retentions (or not) of PMP *w
and *y, yet many investigators do not address them.

There are wide discrepancies of transcription in distinguishing
true vowel-initial words from glottal stop-initial words,
often making comparative work difficult,
and sometimes under-representing \tsc{1sg} or nominal prefixes.
Is it /eu/ ‘I go’? Or /ʔ-eu/ ‘\tsc{1sg}-go’? Is it /ai/
and /au/ ‘tree, wood’? Or /ʔai/
and /ʔau/ (reflecting the historical *k of PMP *kahiw)?
And how do we know? \citet[60--62]{Tamelan2021} provides several tests
for the \ili{Dela} language to determine whether a glottal is phonemic or not.
There are some languages in which PMP *k > \it{ʔ},
but it might not be written in any position by the investigators,
or is written only between vowels but not initially or finally,
or is written only irregularly although it actually occurs regularly.
We have encountered all of these issues.
There are other languages in which PMP *k > {\0} in all positions.
But the distinction between *k > \it{ʔ}
and *k > {\0} can be significant for comparative purposes and for subgrouping.
These examples also illustrate that orthographies used for linguistic
resources should perhaps be different on some points from practical
orthographies for popular consumption.

Several older Dutch and German sources in the region indicate
an intervocalic glottal with a diaeresis over the second vowel.
So orthographic ‹seï› for /seʔi/,
‹haü› for /haʔu/, ‹foeöe› for /fuʔu/,
‹djaä› for /ʤaʔa/, etc. (see \citealp{Grimes1999}).
Yet in other sources the diaeresis may represent that the
two vowels are articulated as separate syllables,
rather than coalescing as a phonetic diphthong.

Given so much variation in transcription,
consequently in several Holle lists it is not surprising that
the introductory notes indicate a diaeresis over a vowel
indicates an intervocalic glottal.
Yet this is also said for languages that have no intervocalic
glottals (such as \ili{Buru}
and others in central Maluku).
Diacritics of various sorts (with little consistency across languages),
are used to indicate: glottal,
stressed syllable, unstressed syllable,
length, schwa, close vowel, open vowel,
cancelling Dutch orthographic practices (i.e.
a diacritic may indicate ‹oe› should be read as a vowel
sequence /oe/, not as /u/), and other things.
In different lists (and sometimes in the same list),
‹a›, ‹á›, ‹à›, ‹ä›,
‹â›, ‹å›, ‹ā›, ‹a›, and ‹aa›,
could indicate the same thing,
or they could indicate different things.
One must know something about the language to decode the diacritics,
and phonemic differences need to be retained in the documentation
and reflected in comparative-historical discussion.

Many Austronesian languages in the target region have VV sequences of like vowels.
These are phonetically longer than a single V,
and measured instrumentally average around 20--45{\%} longer
in duration than a single V
(see \citealp[95--99]{Edwards2020}
and \citealp{Grimes2024}).
The duration of VV sequences of like
and unlike vowels is often only trivially different,
with unlike sequences tending to be slightly longer.
Phonemically these double VV are often analysed as two syllable
peaks, and this is reflected in the patterns of word stress,
affixation, compounding, and reduplication \citep{CulhaneEdwards2021}.\footnote{
		Unfortunately this very common
		and significant feature of the vowel systems of the region is
		only mentioned in passing by \citet{Hajek2010}.}

Native speakers may be inconsistent in writing the difference
between double vowels and single vowels, as their main orthographic model is usually
Indonesian -- a language which does not distinguish double vowels.
However, when they become aware of the difference,
they consistently write the differences between a single
and double vowel in many languages in the region.
Some same VV sequences arise from the loss of a historical
consonant as illustrated in \trf{02-07}
(STB = \tsc{Seram-Tanimbar-Bomberai}).

\begin{table}
	\centering\tcp{VV sequences from loss of historical *C (consonant)}{02-07}
	\st{4.25pt}\begin{tabular}{ll|l@{ }l|l@{ }l|l}\lsptoprule
&&&&&&trunk, base,\\
&local gloss&\mc{2}{l|}{sail}&\mc{2}{l|}{name}&source\\
Subgroup&PMP&\mc{2}{l|}{*layaR}&\mc{2}{l|}{*ŋajan}&*puqun\\ \midrule
\tsc{\ili{Bima}-Lembata}&\ili{Dhao}&\mc{2}{l|}{{\hr}lai}&\mc{2}{l|}{{\hr}ŋara}&\tcg{(kapua)}\\
\tsc{Timor-Babar}&\ili{Tii}&\mc{2}{l|}{{\hr}laa}&\mc{2}{l|}{{\hr}nade}&{\hr}huu\\
\tsc{Timor-Babar}&\ili{Tetun}&\mc{2}{l|}{{\hr}laa-n}&\mc{2}{l|}{{\hr}nara-n}&{\hr}huun\\
\tsc{\ili{Central} Timor}&\ili{Kemak}&\mc{2}{l|}{}&\mc{2}{l|}{{\hr}galan}&{\hr}puun\\
\tsc{Sula-Buru}&\ili{Buru}&{\hr}laa &‘sail’&{\hr}ŋaa-n &‘name’&{\hr}puun\\
&&{\hr}laa-n&‘dorsal fin’&{\hr}ŋaa-t&‘rank’&\\
STB&\ili{Fordata}&\mc{2}{l|}{{\hr}laar}&\mc{2}{l|}{{\hr}nara-n}&\\
STB&\ili{Geser}&\mc{2}{l|}{{\hr}laar}&\mc{2}{l|}{{\hr}ŋasa-n}&\\
		\lspbottomrule
	\end{tabular}
\end{table}

It is also the case that long vowels (or disyllabic double vowels)
can be indicated in the primary source material,
but not preserved in the comparative material.
For example, \citet{Jonker1908,Jonker1915}
indicated some long vowels with an acute accent,
others with a grave accent,
and still others not at all.
But even where he did indicate them,
this important information may not be preserved in comparative
works that use him as a source
(see discussion in \citeauthor{Edwards2018a} [\citeyear[366]{Edwards2018a}]
and \citeauthor{Edwards2021} [\citeyear[8--12]{Edwards2021}]).

\citet{BlustTrussel2020} under PCEMP *maya ‘tongue’ list ``Roti''
‹ma-k› and ``Atoni'' ‹ma-f›.
Our own notes, as well as published native-authored materials
for several varieties of Rote,
and several varieties of \ili{Meto} have \it{maa-k}
and \it{maa-f} respectively
(and the same root with other genitive enclitics as well).
*y is regularly lost,
but both historical vowels are retained.
How these words are transcribed can help reveal the historical
patterns, or mask them and contribute to the confusion.

\citet[19-21]{vanKlinken1999} indicates double vowels in \ili{Tetun},
which she analysed phonemically as two syllables,
with an acute accent over the vowel,
as in ‹hare› ‘rice’,
‹haré› ‘see’, which is analysed phonemically as /haree/.
For comparative purposes these need to be written as VV sequences of like vowels.
For example, orthographic ‹tún› ‘descend’, needs to be reinterpreted
as \it{tuun} in the comparative literature to track its relationship
with PMP *tuRun ‘descend’,
with the regular loss of *R,
but no loss of a historical vowel. Data sources vary.
When Grimes was driving from Dili to \ili{Kupang} a road sign read
‹hatun velocidade› intending /ha-tuun/ ‘cause to descend = decrease (speed)’.
Without knowledge of the language ‹hatun› could be misinterpreted
as /hatun/, which looks like a reflex of PMP *batu ‘stone, rock’.

In the \ili{Asilulu} dictionary,
\citet{Collins2003a} explains “The colon (:) marks long vowels,
which contrast phonemically with short vowels.” (p.
xxvii) “perhaps a more consistent orthography would have represented
these long vowels as double vowels (\it{aa},
\it{ii}, and so forth).” (p. xxxvi).
It is also the case that for a single language,
one investigator may indicate the long/double vowels,
while another investigator may be oblivious to their existence.

Informed clarifications are extremely helpful,
such as the information found in \citet{Nivens2017} that in the
comparative \ili{Aru} data, transcriptions are adapted to IPA standards,
with the following exceptions: the letter ‹y› is used for
the palatal glide [j],
\tsc{Tarangan} mid-close [e] and [o] are transcribed ‹ê›
and ‹ô›; mid-open [ɛ]
and [ɔ] are transcribed ‹e›
and ‹o›; and unexpected stress is marked by an acute accent (the
normal pattern is penultimate stress,
or final stress if the final syllable contains a mid-vowel).
The symbol ‹ʍ› represents a voiceless labiovelar fricative.
Such clarification further illustrates our challenges for unravelling
the wide variations of what is indicated throughout the region
by various diacritics over vowels.
Unfortunately, many sources of data provide no information about
the transcription practices employed.

In unravelling all this confusion,
we have often had to triangulate with multiple sources,
including consulting more modern sources,
more reliable sources on neighbouring languages,
consulting with scholars who have worked on those languages to
clarify what has been said by them or by others,
and occasionally we have found native speakers or gone to the
region to better understand the phenomena hinted at in the literature
or transcription of those languages.



	\sscp{Issues with languages that have more than five vowels}{02-04-06}
Across the \tsc{\ili{Bima}-Lembata} subgroup of Austronesian languages,
and also in the Papuan \tsc{\ili{Alor}-Pantar} languages,
in addition to phonemic contrasts between phonetically short
(single) and long (double) vowels,
there are other widespread patterns that involve schwa /ə/,
reduced vowels (of shortened duration),
centralised vowels (where the five vowels move slightly towards
schwa /ə/), de-stressing of penultimate syllables with schwa /ə/,
or sometimes a stressed penultimate schwa /ə/,
and phonetic lengthening of the following intervocalic consonant.
In several languages there is an interaction between several of these phenomena.

There are a variety of strategies found in various sources
for the transcription of these different and very subtle phenomena.
This is further complicated as Indonesian has a phonemic
six-vowel system with schwa /ə/,
but /ə/ and /e/ are orthographically
under-differentiated and both are represented by ‹e›
\citep{Grimes1996}.\footnote{
		The mid-front vowel /e/ has a very low functional load in Indonesian
		(low frequency), whereas /ə/ has a very high functional load.}
Many native speakers
of local languages, or Indonesian scholars
trying to document languages
in this region often transcribe their data with only the five
vowels of the typewriter or basic keyboard,
and no other indication.

In languages with schwa /ə/,
reduced vowels, or centralised vowels, sometimes nothing is indicated,
even though there is phonemic contrast in identical environments.
Sometimes the investigator marks something on the vowels.
Such marking may or may not be consistent.
Sometimes the investigator writes double consonants to indicate
the phonetic lengthening of the following consonant which is
often triggered by a preceding schwa /ə/,
and is entirely predictable.
However the schwa that triggers the phonetic lengthening is not
indicated by many attempting to write or document those languages.
Such a strategy is somewhat effective for words shaped CəCV → ‹CaCCV›
or ‹CeCCV›, but it does not work for words shaped CəV or əV which have no
following consonant to double.

The \tsc{\ili{Bima}-Lembata} subgroup is sandwiched between other Austronesian
languages to the west
and north, and Papuan \tsc{\ili{Alor}-Pantar} languages to the \ili{east}.
\tsc{PAn} is reconstructed as having a four-vowel system with
*i, *ə, *a, *u \citep[173]{Blust2013}.
\citet{Smith2021} has pointed out that in many Austronesian languages
in the Philippines, western Indonesia, and Sulawesi penultimate schwa /ə/ is oftenː
a) unstressed, b) of short duration,
and/or c) accompanied by phonetic lengthening
of the following consonant.

Besides the patterned differences between the phonetic realisation
of schwa in \tsc{\ili{Bima}-Lembata} languages,
it may be significant that the phonetic realisation of schwa
in these languages is very similar to vocalic phenomena
in nearby Papuan \tsc{\ili{Alor}-Pantar} languages such as \ili{Klon}, \ili{Wersing}, and \ili{Klamu}.

The issue with the transcription of schwa
and reduced vowels in \tsc{\ili{Alor}-Pantar} languages is
seen with alternate language names such as \ili{Klon}, Kelong, Kalong.
\citet[35]{Baird2008} notes that in \ili{Klon} “schwa only occurs in unstressed
syllables.” At a phonology-orthography workshop in November 2019
with a number of native speakers of the \ili{Klon} \ili{Bring} dialect,
it was found that \ili{Klon} has seven phonemic vowels /i ɪ e a o ʊ
u/ with only one central vowel /a/
and contrastive vowel length in all seven vowels.
There is no phonemic schwa,
but there is a non-phonemic epenthetic vowel used to break up consonant clusters.
It is unstressed, of short duration and not syllabic.
It is mostly central,
and often takes on the colour of nearby vowels \citep{Jones2020a}.
Native speakers sometimes refer to
this as a “fast vowel” (Indonesian: \it{vokal cəpat}).
So /mgol/ ‘banana’ → [məgol] {\tl} [mŏgol] {\tl} [mgol],
/ngan/ ‘thing, matter’ → [nəgan] {\tl} [năgan] {\tl} [ngan],
/ngeen/ ‘give us’ → [nəgeːn] {\tl} [nĕgeːn] {\tl} [ngeːn], etc.

Although the Papuan language of \ili{Wersing} spoken on \ili{Alor} has only
a five-vowel base system with long
and short contrasts, it also has a similar phenomenon of epenthetic
schwa-like “fast vowels” used to break up consonant clusters,
as in /ptori/ → [pətori] ‘boiled corn’,
/t-teŋ/ → [təteŋ] ‘our hands’,
/tma/ → [təma] ‘sea’ \citep{Jones2020b}.

The interplay between vowel quality,
vowel length, and consonant length appears to be an areal phenomenon
found in the \tsc{\ili{Bima}-Lembata} nodes of \tsc{Flores-Lembata},
\tsc{Sumba-Havu}, and also in some Papuan \tsc{\ili{Alor}-Pantar} languages.
In addition to lengthening of the following consonant in \tsc{Sumba-Havu}
following a penultimate schwa described in \srf{03-02-01} and \srf{03-03-01},
the languages of western Sumba show a complex interplay between
the quality and length of the penultimate vowel,
and some phonetic lengthening of the following intervocalic consonant.
In some cases all that remains as traces of a certain historical
vowel is the lengthened consonant (see \srf{03-03-03-01}).

\largerpage
\ili{Klamu}, spoken on \ili{Pantar}, manifests most of the phenomena discussed
above in a single language.\footnote{
		There is some preliminary indication that other languages
		of \ili{Pantar}, such as \ili{Teiwa}, may have similar
		complexities to \ili{Klamu} (Amos Sir, pers.
		comm.).
		The “fast vowel” is found across most \ili{Alor}-\ili{Pantar} languages.}
\ili{Klamu} has regular vowels,
long vowels, and short penultimate vowels that trigger phonetic
lengthening of the following consonant.
The five base vowels thus show a three-way contrast in length,
and there is interplay between the short vowels with compensatory
lengthening of the following consonant.
There is also an epenthetic ``fast vowel'' with initial CC clusters
similar to that described for \ili{Wersing} and \ili{Klon}.
The short vowels are distinguished by three phenomena: a) they
are of slightly shorter duration [ă ĕ ĭ ŏ ŭ] than the equivalent
base vowel (V); b) there is vowel reduction in which the short
V are slightly centralised in vowel quality (moving towards schwa
[ə], but only the low vowel /a/ actually reaches a full,
albeit slightly shortened [ə]; c) this shorter duration
and slight centralisation combines with a slight lengthening
[ˑ] of the following C.\footnote{
		The lengthening of the following
		C in \ili{Klamu} is indicated by a half-length [ˑ],
		and to Grimes' ears is not as long as the lengthening of the
		C [ː] following /ə/ in \ili{Havu}
		and \ili{Dhao} (see discussion
		and examples in \srf{03-03-01-02} and \srf{03-03-02-HD}).}
This combined phenomenon
is indicated in the locally derived orthography by a doubling
of medial consonants (including digraphs),
with no marking on the preceding V.
Doubling of the consonants is also found in the writing of family
names, indicating the orthographic strategy has been around for some time.
Word stress is normally penultimate,
including on the reduced V.
Almost all of the hundreds of examples found involve combinations
of /i, a, u/.
Very few involve mid vowels /e, o/.
These phenomena are illustrated in example \qf{02-44}.
(Data from Grimes' fieldnotes.)

{\wg\begin{xltabular}{\linewidth}{>{\hspace{-\tabcolsep}}l<{\hspace{-\tabcolsep}}>{\hspace{-\tabcolsep}\enspace}llX}
\lsptoprule
\tbx{02-44}		&	Orth.	&	Phonetic	&	Gloss		\\ \midrule
\endhead
&	‹wala›	&	[ˈwɑlɑ]	&	‘rock, stone’		\\
	&	‹walla›	&	[ˈwălˑɑ] {\tl} [ˈwəlˑɑ]	&	‘full, packed, crowded’		\\
	&	‹aci›	&	[ˈɑʧi]	&	‘peel skin (of animal or plant)’		\\
	&	‹acci›	&	[ˈăʧˑi] {\tl} [ˈəʧˑi]	&	‘undo knot’		\\
	&	‹baqa›	&	[ˈbɑqɑ]	&	‘hole, opening, orifice’		\\
	&	‹baqqa›	&	[ˈbăqˑɑ] {\tl} [ˈbəqˑɑ]	&	‘body’		\\
	&	‹raaqu›	&	[ˈrɑːqu]	&	‘pick, harvest (grain)’		\\
	&	‹raqqu›	&	[ˈrăqˑu] {\tl} [ˈrəqˑu]	&	‘two’		\\
	&	‹rana›	&	[ˈrɑnɑ]	&	‘face to face’		\\
	&	‹ranna›	&	[ˈrănˑɑ] {\tl} [ˈrənˑɑ]	&	‘how much, how many’		\\
	&	‹lala›	&	[ˈlɑlɑ]	&	‘promise’		\\
	&	‹lalla›	&	[ˈlălˑɑ] {\tl} [ˈləlˑɑ]	&	‘loose’		\\
	&	‹lilla›	&	[ˈlɪlˑɑ] {\tl} [ˈlĭlˑɑ] {\tl} [ˈləlˑɑ]	&	‘sick, hurt’		\\
	&	‹iila›	&	[ˈiːlɑ]	&	‘water’		\\
	&	‹illa›	&	[ˈɪlˑɑ] {\tl} [ˈĭlˑɑ] {\tl} [ˈəlˑɑ]	&	‘fly (v)’		\\
	&	‹banga›	&	[ˈbɑŋɑ]	&	‘ask, request, bargain’		\\
	&	‹bangnga›	&	[ˈbăŋˑɑ] {\tl} [ˈbəŋˑɑ]	&	‘life’		\\
	&	‹wana›	&	[ˈwɑnɑ]	&	‘connect’		\\
	&	‹wanna›	&	[ˈwănˑɑ] {\tl} [ˈwənˑɑ]	&	‘exists’		\\
	&	‹billi›	&	[ˈbɪlˑi] {\tl} [ˈbəlˑi]	&	‘see, check out, look in on’		\\
	&	‹gagga›	&	[ˈgăgˑɑ] {\tl} [ˈgəgˑɑ]	&	‘he, she, it (one of several \tsc{3sg} pronouns)’		\\
	&	‹ecicci›	&	[ɛˈʧɪʧˑi] {\tl} [ɛˈʧəʧˑi]	&	‘true’		\\
	&	‹akku›	&	[ˈăkˑu] {\tl} [ˈəkˑu]	&	‘say, that (complementiser for speech acts)’		\\
	&	‹cukku›	&	[ˈʧŭkˑu] {\tl} [ˈʧəkˑu]	&	‘enter, go into’		\\
	&	‹gsarru›	&	[ˈgsărˑu] {\tl} [ˈgsərˑu]	&	‘see-him (undergoer)’ /g-sarru/		\\ \hhline{~}
	&						&	{\tl} [gəˈsərˑu]	&	\hp{ }		\\
	&	‹tlaqqu›	&	[ˈtlăqˑu] {\tl} [ˈtləqˑu]	&	‘fever and chills,		\\ \hhline{~}
	&						&	{\tl} [təˈləqˑu]		&		\hp{`}symptoms of malaria’	\\
	&	‹dumma›	&	[ˈdŭmˑɑ] 	&	‘much, very, great		\\ \hhline{~}
	&	\hp{ }		&	\hp{ }	&	\hp{`}(amount of, degree of)’		\\
	&	‹yedda›	&	[ˈyĕdˑɑ] {\tl} [ˈyədˑɑ]	&	‘1) still; 2) not yet’		\\
	&	‹afenni›	&	[ɑˈfĕnˑi] {\tl} [ɑˈfənˑi]	&	‘arrogance’		\\
\lspbottomrule
\end{xltabular}\addtocounter{table}{-1}}

Here we simply note the close phonetic similarities in the description
of \tsc{\ili{Alor}-Pantar} languages with descriptions of schwa in
nearby Austronesian languages on Flores as being unstressed,
of short duration, like an epenthetic vowel to break up consonant
clusters, and sometimes triggering lengthening of the following consonant.
We are not claiming contact-induced change in either direction,
but that is something for further investigation.
Since this schwa of short duration often triggering lengthening
of the following consonant was already found in Austronesian
languages to the west
and north, at best its presence in Papuan languages lends support
to maintain it in some Austronesian languages within a limited region.
The phenomena involving schwa in \tsc{\ili{Bima}-Lembata} languages are discussed
in \srf{03-02-01} and \srf{03-03-01}.

%In \srf{02-04-06-01} we discuss various ways people indicate
%schwa /ə/ or other vowels in writing languages with six-vowel
%systems in the region.
%In \srf{02-04-06-02} we discuss particular issues encountered
%in the documentation of Sumba languages,
%along with other transcription issues complicating understanding
%languages in the \tsc{\ili{Bima}-Lembata} subgroup.
%In \srf{02-04-06-03} we discuss particular issues for understanding
%\citeauthor{Onvlee1984}'s (\citeyear{Onvlee1984}) transcriptions,
%necessary to make comparative use of his valuable data source.
	\sscp{Different sound correspondences in language/dialect clusters}{02-04-07}
Different varieties of larger “languages” or clusters of closely
related languages and dialects often have different historical sound correspondences.
This is true, for example,
for the languages of Rote Island
and the \tsc{Meto} languages of western Timor.
Many of these systematic differences have been known for
some time \citep{Jonker1908,Jonker1915}.
These systematic distinctions are very important for unravelling
the prehistory of the region \citep{Edwards2018a,Edwards2018b,Edwards2021}.
Yet comparative sources that simply say “Rote” (or “Roti” [sic])
or “Uab \ili{Meto}” (or “Dawan” or “Atoni’”),
and fail to specify which variety,
may mask important differences,
rather than track systematic patterns rigorously.
For example, PMP *\tbr{z}a\tbb{l}an ‘path, way’ >
Roi{\Q}is \ili{Amarasi} \it{\tbr{r}a\tbb{n}an},
%\ili{Kotos} \ili{Amarasi} \it{\tbr{r}a\tbb{n}an},
\ili{Amanuban} \it{\tbr{l}a\tbb{n}an},
and Amfo{\Q}an \it{\tbr{l}a\tbb{l}an} reflecting different correspondences
of PMP *z and *l.
In the past these have often been treated as a single language
in spite of cumulatively significant differences in the sound
correspondences, lexicon, morphophonemic processes, and grammatical functors.
The \tsc{\ili{Ambon}-Seram} language \ili{Alune} spoken on Seram also has
considerable variation in its many varieties
and in their historical sound correspondences.
\scp{Issues in methodology}{02-05}
An assumption of a monolithic spread of the Austronesians into
the region can easily make it difficult to see the discontinuities
that are evident in the data.
An assumption of a single higher level subgroup encompassing
most or all of the Austronesian languages of our target region
can make it hard to see that sporadically similar patterns are
not actually exclusively shared innovations,
once the sound correspondences are mapped out
and fuller data are examined.

Conversely, an assumption that there is not a single encompassing
subgroup in a region can make it difficult to see similarities
and links that are actually there.
An assumption that lower level subgrouping arguments posited
by others “must be correct”,
needs to be held with caution.
We ourselves found several cases in which,
out of respect for the work of others,
we initially adopted their subgroups
and worked with those.
But in following where the data led us,
(often with additional, more complete
and more reliable data),
we found there were other issues that had to be given greater
weight, and eventually we had to depart from earlier classifications,
redefine the Wallacean subgroups (or nodes within them),
and even rewrite and reorganise chapters already drafted in this book.

We also need to be cautious about jumping to conclusions about
time depth based on the linguistic evidence.
While certain sound changes can inform us about 
a relative sequence of which ones 
came first and which ones came later
(e.g.
*b > \it{f > h},
*p > \it{ɸ > h >} {\0},
*k > \it{ʔ} > {\0},
*s > \it{h}) on their own they cannot tell us at
what point in time these changes happened.
%While certain sound changes can tell us about a relative sequence
%of what most likely was first,
%and what came later (e.g.
%*b > \it{f > h},
%*p > \it{ɸ > h >} {\0},
%*k > \it{ʔ} > {\0},
%*s > \it{h}), on their own they cannot tell us when that happened.
The field of glottochronology has been largely abandoned because
we know, for example, that a) languages do not change at a constant rate
through time, b) different languages change at different rates,
and c) there can be spurts of rapid change,
for example, associated with war
and conquest or tragedy such in the aftermath of volcanic eruptions and tsunamis.
But sometimes a datable change is extrapolated to other languages,
and this may or may not be a problem.
For example, earlier it was noted that \citet[10]{Klamer2017} linked
the break-up of the \tsc{\ili{Alor}-Pantar} Papuan languages to around
the same time as the arrival of the Austronesians because “the
reconstructed vocabulary of proto-AP appears to contain Austronesian
loan words such as ‘betel nut’ [{\ldots}] following regular sound
changes suggest that the AP family split up after being in contact
with the Austronesian languages in the area.”

That is certainly possible,
but \citet[81]{Grimes1991c} has also noted that several Portuguese
loans in \ili{Buru} display the same changes to sounds
and word structure that are found in PMP words.
So Portuguese \it{meirinho} /maiɾiɲu/ ‘bailiff’ > \ili{Ambon} \it{marinyu}
> \ili{Buru} \it{emrimo} ‘village crier, constable’.
And Portuguese \it{capitão} /kapitãũ/ ‘ship captain,
military officer’ > \ili{Ambon} \it{kapitan} > \ili{Buru} \it{epkitan} ‘fighting champion, wartime leader’.
Similar restructuring is seen in \ili{Buru} loans ultimately from \ili{Sanskrit},
such as: \it{baksa} ‘read’,
\it{eʔwasa-t} ‘wealth, power, influence’, \it{fisara} ‘speak’
(\ili{Lisela} \it{epsara}), \it{mansari} ‘hunt’, \it{mansia} ‘human’,
and \it{storita} ‘story, tell a story’.
And also in \ili{Buru} loans ultimately from \ili{Arabic},
such as: \it{barkate} ‘bless (v),
blessing (n)’, \it{eʔʤadi} ‘become,
be born’, \it{eʔrama-t} ‘devout,
spiritual’, \it{eslau} ‘beg forgiveness, restore right relationships’,
\it{haiwane} ‘domestic animals’, \it{hasa-k} ‘oppose,
hate’, \it{hum sikit} ‘mosque’, \it{ʤunai} ‘world’,
\it{kufur} ‘grave’, \it{slame} ‘Islam,
Muslim’, \it{srane} ‘Christian, christen’.

\citet[397]{Edwards2021} could reconstruct a Proto-Rote-\ili{Meto} form **ŋgarei ‘church’
based on regular sound correspondences.
Yet he has not done so,
since he is aware that the daughter forms are from Portuguese \it{igreja} [iˈgreʒa].
Portuguese contact did not happen before 1512 AD,
and \ili{Sanskrit} and \ili{Arabic} influence came significantly later than
Austronesian incursion, while the connection to PMP in the region goes back about 4,000 years.
So parallel developments or processes in the phonotactics
and historical phonology may not necessarily reflect similarities in time depth.

Papuan-sourced {\#}muku ‘banana’ and Austronesian-sourced reconstructions
such as ‘betelnut’ in Papuan TAP languages provide evidence that
words can and do spread across unrelated groups that are not sourced in
a common parent language or linkage.
Additional similar examples are given in the later chapters of this book.
In a geographically diverse
and dispersed region with known trade
and travel, these examples should encourage serious caution (and
scepticism) for using lexical evidence as the basis for subgrouping arguments.


\largerpage
\scp{“Eastern Indonesia” as a complex ethnographic transition zone}{02-06}
Following on from the discussion of contact-induced change in
the Austronesian languages of the region (\srf{02-03}),
it is important to note that evidence of contact is not limited
to language, but is also found in a range of distinct socio-cultural features
that have been observed for nearly 200 years \citep{Fox2015}.
\citet[1]{vanWouden1935} defined Eastern Indonesia as a distinct
ethnographic region “based upon various points of similarity
from an ethnographical point of view”.
The similarities he noted included a clan system,
exclusive cross-cousin marriage, and sociocosmic dualism,
collectively being features that set these societies apart from
Austronesian societies to the west.
\citet{Fox1980,Fox1989,Fox2015} revisited van Wouden’s work,
expanding the list of distinguishing features to include:

\begin{itemize}
	\item Concern with origins;
	\item Use of dual categories of complementary opposition;
	\item Identification of life processes and relationships in botanical metaphors; (also \citealp{Fox1971});
	\item The notion of ‘person’ as composed of \it{blood / flesh} and \it{semen / bone};
	\item Diarchy predicated on delegation or usurpation of authority;
	\item “House” as a primary descent group;
	\item Clans form political or ceremonial groupings;
	\item Marriage alliances link houses/clans;
	\item Alliances are viewed as the “flow of life”;
	\item Extended series of life-enhancing rituals joining predecessor to descendant in a cyclical translation of life.
\end{itemize}

Looking eastward into Melanesia
and south into northern Australia,
\citet[18]{Grimes2010} concludes “Eastern Indonesia is a zone where
social interaction has resulted in shared cultural heritages
with other Austronesian societies as well as non-Austronesian
societies.” Her conclusion for \ili{Buru} society (in central Maluku)
is that it could be described as an Austronesian speaking society
with a significant number of features common to societies speaking
Papuan languages in the greater Melanesian region,
including in the highlands of Papua New Guinea
and islands \ili{east} of New Guinea.
Whereas \ili{Fataluku} (in Timor-Leste) is “a society speaking a Papuan
language with many cultural features in common with their Austronesian
speaking neighbors \citep{McWilliam2007}.” From an ethnographic perspective,
there are additional examples supporting the view that eastern
Indonesia should be considered as a complex zone of transition \citep[191]{Fox2015}.
While some of these features by themselves are found in many
parts of the world,
this unique bundling of cultural features is not widespread in
Austronesian societies west of the Blust line.

More specifically, \citet{Fox2015} finds that the distribution
in terminologies and semantics relating to opposite sex siblings
and to the term for `child’s spouse’s parent' align roughly with the Blust line.
These are touched on in more detail in subsequent chapters.

While *mamah ‘father’s brother’ is reconstructed to \tsc{PAn},
in Wallacea many languages have developed a lexicalised term
for ‘mother’s brother’, not found in the west.
In some Wallacean languages the term is cognate with
\citeauthor{Stresemann1927}'s (\citeyear{Stresemann1927}) PAMS **məmə ‘mother’s brother’.
Other terms for ‘mother’s brother’ are also found.
This is an example of a \it{category innovation} (not a lexical
innovation with formal cognates),
that is likely to arise from substrate influence with variation in different regions.
This concept is actually much more nuanced.
For example, \ili{Buru} (central Maluku) \it{meme} ‘mother’s brother’
is a reflex of PAMS **məmə ‘mother’s brother’.
\ili{Tii} (Rote Island) \it{toʔo} ‘mother’s brother’ is not (it is
from *tau ‘person’ [\citealp[370]{Edwards2021}]),
but captures exactly the same concept or category in the kinship system.
Furthermore both languages use the same botanical metaphor from
‘tree, base of tree, source’ to point to the ‘mother’s brother with
the most senior rank, eldest’.
So \ili{Buru} \it{meme lahi {\tl} mem-lahin} ‘mother’s oldest (classificatory)
brother (lit: root/source uncle)’ (\it{lahin} ‘root; tree,
plant; base, source, cause’), is not formally cognate with,
but has almost identical semantics to,
\ili{Tii} \it{toʔo huu-k} ‘mother’s oldest (classificatory) brother
(lit: root/source uncle)’ (\it{huu-k} ‘root; tree,
plant; base, source, cause’; cf.
verbal doublet \it{ihuu-ioka} ‘take root’).
The category of ‘mother’s brother’ is widespread,
even though only some of the terms are formally cognate.
For example, \ili{Amarasi} \it{baba-f} ‘mother’s brother,
father-in-law’, \ili{Helong} \it{baki} ‘mother’s brother’,
\ili{Kemak} \it{naa-r} ‘mother’s brother (from *ñaRa ‘brother (f.s.))’,
\ili{Raklungu} \it{mimi} ‘mother’s brother’,
\ili{Selaru} \it{memi} ‘mother’s brother,
father-in-law’ reflect a very stable concept across a number of subgroups.

Further evidence is found in patterns of musical forms.
\citet{Blench2021} notes that polyphony is found throughout the
Austronesian cultures of Formosa (except one),
but was apparently left behind there when the Austronesian speaking
people moved out of Formosa.
However, vocal and instrumental polyphony is also widespread
where Papuan languages are spoken.
Referring to forms of polyphony found in Austronesian speaking
societies in eastern Flores
and parts of Timor,
Blench finds it “plausible that patterns found in eastern Indonesia
reflect the Papuan substrate”.

There is an abundance of direct evidence in regions where Papuan
languages are still spoken today,
and an abundance of trace evidence where there are no Papuan
languages spoken today, that allow us to posit contact as the source or trigger for a
number of features in the lexicons,
grammars, and even in the cultures of Austronesian-speaking peoples.
These features are found in the Papuan languages.
They are not generally found in the Austronesian languages west
of the Blust line.
The nature of the evidence implies intense contact,
bilingualism and linguistic and cultural accommodation over a
long period of time.
However they do not address the issues of whether or not there
was a shared parent language,
nor to subgrouping issues in the Austronesian languages in the region.

\scp{A mixed metaphor}{02-07}
The issues above reflect only a sample of things that have added
confusion rather than clarity in unravelling the historical-comparative
picture of the region.
Most of the arguments about linguistic prehistory
and subgrouping in the region to date have been based on very little data.

In looking at Austronesian speaking peoples speaking mostly typologically
SVO languages coming into a relatively large
and diverse region already populated by people speaking various
typologically SOV Papuan languages,
we have a picture something like a river of fresh water coming
into contact with a body of salt water (see \frf{02-28}).
There are places where the salt water is demarcated
and is clearly salt water,
and the fresh water is clearly fresh water.
And places where it is something different altogether -- it is brackish.
The fresh water is not uniform where it comes into the sea,
with a deep channel on the side where the water is running swiftly,
shallower and running more slowly on the opposite side,
with multiple tributaries at the mouth flowing around small islands,
large rocks and sandbars.
The sea is also not uniform with different temperatures
and ecosystems at different distances from shore
and at different depths.

\begin{figure}
	\centering
	\includegraphics[width=0.45\columnwidth]{02-28.jpg}
	\fcp{Wae Tina Nangan river mouth in south Buru}{02-28}
\end{figure}

Once the waters meet,
there are places where the identity as one or the other is not
always clear, because they have mixed; they are no longer purely fresh water
or purely sea water.
They start looking a bit the same in places.
But the mixture is not the same throughout the zone of contact.
Different bottom features (substrate),
channels, islands, eddys and back currents mean that there is
great variety within the area where the two meet.
And silt and debris floating around or stirred up by the two
currents coming together add distractions
and dangers to the unwary.
And there is not just one river flowing into the sea.
For example, not far to the \ili{east} of the river mouth shown in
\frf{02-28} is another small river,
Namrole (Role Bay -- that gives the large bay its name),
flowing into the same sea.
And not far to the west of the picture is Wamkana (Forceful River),
a small river flowing with great force.

Experienced boatmen know they must be extra careful where rivers
flow into the sea,
because of the unpatterned waves
and currents, and unexpected debris that can flip or damage a boat.
They also know that the intersection of the two bodies of flowing
water means the hidden bottom is ever-changing
and a bit unpredictable -- not the same as the last time they passed by.
